\subsection{H. WEYL 积分公式}

令 e 为 G 的单位元，\(\bar{e}\) 为其在 G/T 中的类。将 G 在 e 处的切空间与 g 识别，将 T 在 e 处的切空间与 t 识别，并将 G/T 在 \(\bar{e}\) 处的切空间与 g/t 识别。将 \((u,v)\mapsto u\cap v\) 记为下列 R-双线性映射

\[
\operatorname{Alt} ^ {n - r} (\mathfrak {g} / \mathfrak {t}) \times \operatorname{Alt} ^ {r} (\mathfrak {t}) \rightarrow \operatorname{Alt} ^ {n} (\mathfrak {g})
\]

它在第 1 节中已有定义。

回忆（第 III 章，第 3 节，第 13 号，命题 50），映射 \(\omega \mapsto \omega(e)\) 是从 G（分别为 T）上的 \(n\)（分别为 \(r\)）次左不变微分形式的向量空间到 \(\mathrm{Alt}^n(\mathfrak{g})\)（分别为 \(\mathrm{Alt}^r(t)\)）的同构。进一步注意，由于 \(\mathbf{R}^*\) 的每个连通紧子群都退化为单位元，\(\det \mathrm{Ad}g = 1\) 对所有 \(g \in \mathrm{G}\) 成立，所以 G 上的 \(n\) 次左不变微分形式同时也是右不变的，并且在内自同构下不变（第 III 章，第 3 节，第 16 号，命题 54 的推论）：从现在起我们简称为 G-不变微分形式。

同样，由第 III 章第 3 节第 16 号命题 56 及上述结果可知，映射 \(\omega \mapsto \omega (\bar{e})\) 是从次数为 \(n - r\) 的 \(\mathrm{G / T}\) 上的 G-不变微分形式空间到 \(\mathrm{Alt}^{n - r}(\mathfrak{g} / \mathfrak{t})\) 的同构。

若 \(\omega_{G/T}\) 是 G/T 上次数为 n - r 的 G-不变微分形式，而 \(\omega_{T}\) 是 T 上次数为 r 的不变微分形式，则将 G 上满足下式的唯一 n 次不变微分形式记为 \(\omega_{G/T} \cap \omega_{T}\)：

\[
(\omega_ {\mathrm{G/T}} \cap \omega_ {\mathrm{T}}) (e) = \omega_ {\mathrm{G/T}} (\bar {e}) \cap \omega_ {\mathrm{T}} (e).
\]

最后回忆，\(f : (\mathrm{G}/\mathrm{T}) \times \mathrm{T} \to \mathrm{G}\) 表示如下流形态射：它由映射 \((g, t) \mapsto gtg^{-1}\) 从 \(G \times T\) 到 G 通过取商诱导而来（第 5 节，第 4 号）。若 \(\alpha\) 和 \(\beta\) 分别是 G/T 和 T 上的微分形式，则将 \(\alpha \wedge \beta\) 简记为形式 \(pr_{1}^{*}\alpha \wedge pr_{2}^{*}\beta\) 在 \((\mathrm{G}/\mathrm{T}) \times \mathrm{T}\) 上的形式。

对于 \(t \in \mathrm{T}\)，将 \(\operatorname{Ad}_{\mathfrak{g} / \mathfrak{t}}(t)\) 记为 \(\mathfrak{g} / \mathfrak{t}\) 的自同态，它由 \(\operatorname{Ad} t\) 通过取商诱导而来。置

\[
\delta_ {\mathrm{G}} (t) = \det (\mathrm{Ad} _ {\mathfrak {g} / \mathfrak {t}} (t) - 1) = \prod_ {\alpha \in \mathrm{R} (\mathrm{G}, \mathrm{T})} (t ^ {\alpha} - 1).\tag{3}
\]

令 \(x \in \mathfrak{t}\) 且 \(\alpha \in \mathrm{R}(\mathrm{G},\mathrm{T})\)；将 \(\hat{\alpha}\) 记作元素 \((2\pi i)^{-1}\delta (\alpha)\)（它属于 \(\mathfrak{t}^*\)），于是

\[
((\exp x) ^ {\alpha} - 1) ((\exp x) ^ {- \alpha} - 1) = (e ^ {2 \pi i \hat {\alpha} (x)} - 1) (e ^ {- 2 \pi i \hat {\alpha} (x)} - 1) = 4 \sin^ {2} \pi \hat {\alpha} (x).
\]

若 \(\mathrm{R}_{+}(\mathrm{G},\mathrm{T})\) 表示相对于基 B 的 \(\mathrm{R}(\mathrm{G},\mathrm{T})\) 的正根集合，则有

\[
\delta_ {\mathrm{G}} (\exp x) = \prod_ {\alpha \in \mathrm{R} _ {+} (\mathrm{G}, \mathrm{T})} 4 \sin^ {2} \pi \hat {\alpha} (x),
\]

所以特别地，\(\delta_{\mathrm{G}}(t)>0\) 对所有 \(t\in T_{r}\) 成立。我们还注意到，\(\delta_{\mathrm{G}}(t)=\delta_{\mathrm{G}}(t^{-1})\) 对 \(t\in T\) 成立。

命题 1. 设 \(\omega_{\mathrm{G}}, \omega_{\mathrm{G/T}}\) 和 \(\omega_{\mathrm{T}}\) 分别是 G、G/T 和 T 上次数分别为 \(n, n - r\) 和 \(r\) 的不变微分形式。若 \(\omega_{\mathrm{G}} = \omega_{\mathrm{G/T}} \cap \omega_{\mathrm{T}}\)，则

\[
f ^ {*} (\omega_ {\mathrm{G}}) = \omega_ {\mathrm{G/T}} \wedge \delta_ {\mathrm{G}} \omega_ {\mathrm{T}}.
\]

显然可以假设 \(\omega_{G/T}\) 和 \(\omega_{T}\) 均非零；于是微分形式 \((u,t)\mapsto\omega_{\mathrm{G/T}}(u)\wedge\omega_{\mathrm{T}}(t)\) 在 \((\mathrm{G/T})\times\mathrm{T}\) 上是 n 次且处处非零的；因此存在 \(\delta\) 在 \((\mathrm{G/T})\times\mathrm{T}\) 上的数值函数，使得

\[
f ^ {*} (\omega_ {\mathrm{G}}) (u, t) = \delta (u, t) \omega_ {\mathrm{G/T}} (u) \wedge \omega_ {\mathrm{T}} (t).
\]

现在注意，对于 \(h \in \mathrm{G}\)、\(u \in \mathrm{G} / \mathrm{T}\)、\(t \in \mathrm{T}\)，有 \(f(h.u,t) = (\operatorname{Int} h)f(u,t)\)；由于 \(\omega_{\mathrm{G}}\) 在内自同构下不变，立即得到 \(\delta(h.u,t) = \delta(u,t)\)，从而 \(\delta(u,t) = \delta(\bar{e},t)\)。

将商映射 \(p:\mathfrak{g}\to \mathfrak{g} / \mathfrak{t}\) 记为 p，并将由下式定义的映射 \(\varphi :\mathfrak{g} / \mathfrak{t}\rightarrow \mathfrak{g}\) 记为：

\[
\varphi (p (X)) = (\operatorname{Ad} t ^ {- 1}) X - X \text { for } X \in \mathfrak {g};
\]

回忆（第 5 节，第 4 号，引理 4），切映射

\[
\mathrm{T} _ {(e, t)} (f): \mathrm{T} _ {e} (\mathrm{G} / \mathrm{T}) \times \mathrm{T} _ {t} (\mathrm{T}) \to \mathrm{T} _ {t} (\mathrm{G})
\]

将 \((z,tH)\) 送到 \(t(\varphi (z) + H)\)，其中 \(z\in \mathfrak{g} / \mathfrak{t},H\in \mathfrak{t}\)。

令 \(z_{1},\ldots,z_{n-r}\) 为 g/t 中的元素，令 \(H_{1},\ldots,H_{r}\) 为 t 中的元素。则

\[
\begin{array}{l} f ^ {*} \omega_ {\mathrm{G}} (\bar {e}, t) (z _ {1}, \ldots , z _ {n - r}, t H _ {1}, \ldots , t H _ {r}) \\ \quad = \omega_ {\mathrm{G}} (t) (t \varphi (z _ {1}), \ldots , t \varphi (z _ {n - r}), t H _ {1}, \ldots , t H _ {r}) \text {by the calculation of T} _ {(\bar {e}, t)} (f) \\ \quad = \omega_ {\mathrm{G}} (e) (\varphi (z _ {1}), \ldots , \varphi (z _ {n - r}), H _ {1}, \ldots , H _ {r}) \text {since} \omega_ {\mathrm{G}} \text {is invariant} \\ \quad = \omega_ {\mathrm{G/T}} (\bar {e}) (p \varphi (z _ {1}), \ldots , p \varphi (z _ {n - r})) \cdot \omega_ {\mathrm{T}} (e) (H _ {1}, \ldots , H _ {r}) \text {(no. 1, Lemma 1)} \\ \quad = \det (p \varphi) \omega_ {\mathrm{G/T}} (\bar {e}) (z _ {1}, \ldots , z _ {n - r}). \omega_ {\mathrm{T}} (e) (H _ {1}, \ldots , H _ {r}) \\ \quad = \delta_ {\mathrm{G}} (t) \omega_ {\mathrm{G/T}} (\bar {e}) (z _ {1}, \ldots , z _ {n - r}). \omega_ {\mathrm{T}} (t) (t H _ {1}, \ldots , t H _ {r}) \\ \qquad \qquad \qquad \qquad \qquad \qquad \text {since} \omega_ {\mathrm{T}} \text {is invariant} \\ \quad = \delta_ {\mathrm{G}} (t) (\omega_ {\mathrm{G/T}} \wedge \omega_ {\mathrm{T}}) (\bar {e}, t) (z _ {1}, \ldots , z _ {n - r}, t H _ {1}, \ldots , t H _ {r}), \end{array}
\]

所以 \(f^{*}\omega_{\mathrm{G}}(\bar{e},t) = \delta_{\mathrm{G}}(t)(\omega_{\mathrm{G / T}}\wedge \omega_{\mathrm{T}})(\bar{e},t)\)；于是 \(\delta (\bar {e},t) = \delta_{\mathrm{G}}(t)\)，命题得证。

分别以形式 \(\omega_{G}, \omega_{T}\) 和 \(\omega_{G/T}\) 所定义的定向赋予流形 G、T 和 G/T。这些形式在 G、T 和 G/T 上定义不变测度（第 III 章，第 3 节，第 16 号，命题 55 和 56），这些测度也记为 \(\omega_{G}, \omega_{T}\) 和 \(\omega_{G/T}\)。

引理 2. 若 \(\omega_{\mathrm{G}} = \omega_{\mathrm{G / T}}\cap \omega_{\mathrm{T}}\)，则

\[
\int_ {\mathrm{G}} \omega_ {\mathrm{G}} = \int_ {\mathrm{G/T}} \omega_ {\mathrm{G/T}}. \int_ {\mathrm{T}} \omega_ {\mathrm{T}}.
\]

将从 G 到 G/T 的典则态射记为 \(\pi\)。令 \(g \in G\)，并令 \(t_{1}, \ldots, t_{n-r}\) 为 \(\mathrm{T}_{\pi(g)}(\mathrm{G}/\mathrm{T})\) 中的元素。将纤维 \(\pi^{-1}(\pi(g)) = g\mathrm{T}\) 通过平移 \(\gamma(g)\) 与 T 识别。关系 \(\omega_{G} = \omega_{G/T} \cap \omega_{T}\) 于是推出下列等式（《微分与解析流形，结果》，11.4.5）：

\[
\omega_ {\mathrm{G} \sqcup} (t _ {1}, \dots , t _ {n - r}) = (\omega_ {\mathrm{G} / \mathrm{T}} (t _ {1}, \dots , t _ {n - r})) \omega_ {\mathrm{T}}.
\]

因此 \(\int_{\pi}\omega_{\mathrm{G}} = \left(\int_{\mathrm{T}}\omega_{\mathrm{T}}\right)\omega_{\mathrm{G / T}}\)（《微分与解析流形，结果》，11.4.6），并且

\[
\int_ {\mathrm{G}} \omega_ {\mathrm{G}} = \int_ {\mathrm{G/T}} \int_ {\pi} \omega_ {\mathrm{G}} = \int_ {\mathrm{T}} \omega_ {\mathrm{T}}. \int_ {\mathrm{G/T}} \omega_ {\mathrm{G/T}}
\]

（《微分与解析流形，结果》，11.4.8）。

引理 3. 在 \((\mathrm{G}/\mathrm{T}) \times \mathrm{T}_{r}\) 上，\(G_{r}\) 上测度 dg 在局部同胚 \(f_{r}\) 下的逆像是测度 \(\mu \otimes \delta_{G} dt\)（积分，第 V 章，第 6 节，第 6 号），其中 \(\mu\) 是 G/T 上总质量为 1 的唯一 G-不变测度。

在 T（分别为 G/T）上选取最高次数的不变微分形式 \(\omega_{T}\)（分别为 \(\omega_{G/T}\)），使得 \(\omega_{T}\)（分别为 \(\omega_{G/T}\)）所定义的测度等于 dt（分别为 \(\mu\)）。置 \(\omega_{G} = \omega_{G/T} \cap \omega_{T}\)。引理 2 蕴含，\(\omega_{G}\) 所定义的测度等于 dg。令 U 为 (G/T) × T \(_{r}\) 的开子集，使得 \(f_{r}\) 将 U 同构地映到 G \(_{r}\) 的开子集 V。令 \(\varphi\) 为 V 中具有紧支集的连续函数；仍以 \(\varphi\) 表示 \(\varphi\) 延拓到 G \(_{r}\) 且在 V 外为零的函数。有

\[
\begin{array}{r l} & {\int_ {\mathrm{V}} \varphi d g = \int_ {\mathrm{V}} \varphi \omega_ {\mathrm{G}} = \int_ {\mathrm{U}} (\varphi \circ f _ {r}) f _ {r} ^ {*} (\omega_ {\mathrm{G}})} \\ & {\qquad = \int_ {\mathrm{U}} (\varphi \circ f _ {r}) \omega_ {\mathrm{G/T}} \wedge \delta_ {\mathrm{G}} \omega_ {\mathrm{T}} \quad (\text {Prop. 1})} \\ & {\qquad = \int_ {\mathrm{U}} (\varphi \circ f _ {r}) d \mu . \delta_ {\mathrm{G}} d t,} \end{array}
\]

引理得证。

定理 1（H. Weyl）。测度 \(dg\)（定义在 \(G\) 上）是映射 \((g, t) \mapsto gtg^{-1}\) 从 \(G \times T\) 到 \(G\) 下测度 \(dg \otimes \frac{1}{w(G)}\delta_G dt\) 的像，其中

\[
\delta_ {\mathrm{G}} (t) = \det (\operatorname{Ad} _ {\mathfrak {g} / \mathfrak {t}} (t) - 1) = \prod_ {\alpha \in \mathrm{R} (\mathrm{G}, \mathrm{T})} (t ^ {\alpha} - 1).
\]

等价地（积分，第 V 章，第 6 节，第 3 号，命题 4），\(dg\) 是映射 \(f: (\mathrm{G} / \mathrm{T}) \times \mathrm{T} \to \mathrm{G}\) 下测度 \(\mu \otimes \frac{1}{w(\mathrm{G})}\delta_{\mathrm{G}}dt\) 的像。

我们证明最后一个断言。由第 5 节第 1 号以及《微分与解析流形，结果》10.1.3 c) 可知，G-G \(_{r}\) 在 G 中是零测集，T-T \(_{r}\) 在 T 中也是零测集。此外，映射 \(f_{r}\) 使 (G/T) × T \(_{r}\) 成为 G \(_{r}\) 的主覆盖，其群为 W（第 5 节，第 4 号，命题 4 b)）。定理现由引理 3 以及积分，第 V 章第 6 节第 6 号命题 11 得出。

推论 1. (i) 令 \(\varphi\) 为 G 上取值于 Banach 空间或 \(\overline{\mathbf{R}}\) 的可积函数。对于几乎所有 \(t \in \mathrm{T}\)，函数 \(g \mapsto \varphi(gtg^{-1})\) 在 G 上关于 \(dg\) 可积。函数 \(t \mapsto \delta_{\mathrm{G}}(t) \int_{\mathrm{G}} \varphi(gtg^{-1}) dg\) 在 T 上可积，并且

\[
\int_ {\mathrm{G}} \varphi (g) d g = \frac {1}{w (\mathrm{G})} \int_ {\mathrm{T}} \left(\int_ {\mathrm{G}} \varphi (g t g ^ {- 1}) d g\right) \delta_ {\mathrm{G}} (t) d t\tag{4}
\]

（“Hermann Weyl 积分公式”。）

\begin{enumerate}
\def\labelenumi{(\roman{enumi})}
\setcounter{enumi}{1}
\tightlist
\item
  令 \(\varphi\) 为 \(\mathbf{G}\) 上的正可测函数。对于几乎所有 \(t \in \mathrm{T}\)，函数 \(g \mapsto \varphi(gtg^{-1})\) 在 \(\mathbf{G}\) 上可测。函数 \(t \mapsto \int_{\mathbf{G}}^{*} \varphi(gtg^{-1}) dg\) 在 \(\mathrm{T}\) 上可测，并且
\end{enumerate}

\[
\int_ {\mathrm{G}} ^ {*} \varphi (g) d g = \frac {1}{w (\mathrm{G})} \int_ {\mathrm{T}} ^ {*} \left(\int_ {\mathrm{G}} ^ {*} \varphi (g t g ^ {- 1}) d g\right) \delta_ {\mathrm{G}} (t) d t.\tag{5}
\]

由于态射 f 是由映射 \((g, t) \mapsto gtg^{-1}\) 从 \(G \times T \to G\) 通过取商诱导的，只需应用《积分》第 V 章第 5 节、第 6 节、第 8 节以及《积分》第 VII 章第 2 节。

推论 2. 令 \(\varphi\) 为 \(\mathbf{G}\) 上的中心函数（即满足 \(\varphi(gh) = \varphi(hg)\) 对所有 \(g\) 和 \(h\) 属于 \(\mathbf{G}\) 成立），其值属于 Banach 空间或 \(\overline{\mathbf{R}}\)。

\begin{enumerate}
\def\labelenumi{\alph{enumi})}
\item
  \(\varphi\) 可测，当且仅当其在 T 上的限制可测。
\item
  \(\varphi\) 可积，当且仅当函数 \((\varphi |\mathrm{T})\delta_{\mathrm{G}}\) 在 \(\mathrm{T}\) 上可积；在这种情况下，我们有
\end{enumerate}

\[
\int_ {\mathrm{G}} \varphi (g) d g = \frac {1}{w (\mathrm{G})} \int_ {\mathrm{T}} \varphi (t) \delta_ {\mathrm{G}} (t) d t.\tag{6}
\]

将 \(p: \mathrm{G} / \mathrm{T} \times \mathrm{T} \to \mathrm{T}\) 记为第二投影。我们有 \(\varphi \circ f = (\varphi | \mathrm{T}) \circ p\)；此外，\(p\) 下测度 \(\mu \otimes \frac{1}{w(\mathrm{G})} \delta_{\mathrm{G}} dt\) 的像为 \(\frac{1}{w(\mathrm{G})} \delta_{\mathrm{G}} dt\)。推论现由上面的定理 1 以及《积分》第 V 章第 6 节第 2 号定理 1 应用于两个真映射 \(f\) 和 \(p\) 后得到。

推论 3. 令 H 为含 T 的 G 的连通闭子群，h 为其李代数，dh 为 H 上总质量为 1 的 Haar 测度。令 \(\varphi\) 为 G 上取值于 Banach 空间或 R 的可积中心函数。于是函数 \(h \mapsto \varphi(h) \det(\mathrm{Ad}_{\mathfrak{g}/\mathfrak{h}}(h) - 1)\) 在 H 上可积且为中心函数，并且

\[
\int_ {\mathrm{G}} \varphi (g) d g = \frac {w (\mathrm{H})}{w (\mathrm{G})} \int_ {\mathrm{H}} \varphi (h) \mathrm{det} (\mathrm{Ad} _ {\mathfrak {g} / \mathfrak {h}} (h) - 1) d h.\tag{7}
\]

事实上，函数 \(h \mapsto \varphi(h) \det(\mathrm{Ad}_{\mathfrak{g}/\mathfrak{h}}(h) - 1)\) 是 H 上的中心函数，其在 T 上的限制为函数 \(t \mapsto \varphi(t) \delta_{\mathrm{G}}(t) \delta_{\mathrm{H}}(t)^{-1}\)。因此，推论由将推论 2 应用于 G 和 H 得到。

备注。1) 若在推论 3 中取 \(\varphi = 1\)，则得到

\[
\int_ {\mathrm{H}} \det (\mathrm{Ad} _ {\mathfrak {g} / \mathfrak {h}} (h) - 1) d h = w (\mathrm{G}) / w (\mathrm{H})\tag{8}
\]

并且特别有

\[
\int_ {\mathrm{T}} \delta_ {\mathrm{G}} (t) d t = w (\mathrm{G}).\tag{9}
\]

\begin{enumerate}
\def\labelenumi{\arabic{enumi})}
\setcounter{enumi}{1}
\tightlist
\item
  令 \(\nu\) 为商空间 T/W 上由下式定义的测度
\end{enumerate}

\[
\int_ {\mathrm{T/W}} \psi (\tau) d \nu (\tau) = \frac {1}{w (\mathrm{G})} \int_ {\mathrm{T}} \psi (\pi (t)) \delta_ {\mathrm{G}} (t) d t,
\]

其中 \(\pi\) 表示 T 到 T/W 的典则投影。推论 2 意味着，同胚 \(T / W \to G / \operatorname{Int}(G)\)（第 2 节，第 5 号，命题 5 的推论 1）

将测度 \(\nu\) 传送为典则投影 \(\mathrm{G} \rightarrow \mathrm{G}/\mathrm{Int}(\mathrm{G})\) 下测度 dg 的像。

\begin{enumerate}
\def\labelenumi{\arabic{enumi})}
\setcounter{enumi}{2}
\tightlist
\item
  假设 G 单连通。令 A 为 t 的一个单房，令 dx 为 t 上满足 \(\int_{A} dx = 1\) 的 Haar 测度。则测度 \(\nu\) 也可由测度 \(\frac{1}{w(G)} \prod_{\alpha \in R_{+}(G,T)} 4 \sin^{2} \pi \hat{\alpha}(x) dx\) 在 \(\overline{A}\) 上通过同胚 \(\overline{A} \rightarrow T/W\) 传送而得到（第 5 节，第 2 号，命题 2 的推论 1）。
\end{enumerate}

例. 取 G 为群 \(\mathbf{SU}(2,\mathbf{C})\)，取 T 为对角矩阵子群（\(\S3\)，第 6 号）；通过选取 t 的基 \(\{iH\}\)（同上引文）将 t 与 R 识别。置 A = 10, \(\pi\)；这是 t 的一个单房。区间 \(\overline{A} = [0, \pi]\) 可与 G 的共轭类空间识别，其中元素 \(\theta\) 属于 \(\overline{A}\)，对应于矩阵 \(\begin{pmatrix} e^{i\theta} & 0 \\ 0 & e^{-i\theta} \end{pmatrix}\) 所属的共轭类。令 \(d\theta\) 为 \([0, \pi]\) 上的 Lebesgue 测度；由上述结果可知，G 上的 Haar 测度在 \(\overline{A}\) 上的像是测度 \(\frac{2}{\pi}\sin^{2}\theta d\theta\)。

\subsection{李代数上的积分}

命题 2. 令 H 为一个（实）维数为 m 的李群，h 为其李代数。令 \(\omega_{H}\) 为 H 上的 m 次右不变微分形式，并令 \(\omega_{h}\) 为 h 上的 m 次平移不变微分形式，它在原点处与 \(\omega_{\mathrm{H}}(e)\) 相同。我们有

\[
(\exp_ {\mathrm{H}}) ^ {*} \omega_ {\mathrm{H}} = \lambda_ {\mathfrak {h}} \omega_ {\mathfrak {h}}\tag{10}
\]

其中 \(\lambda_{\mathfrak{h}}\) 是 \(\mathfrak{h}\) 上的 Ad(H)-不变函数，满足

\[
\lambda_ {\mathfrak {h}} (x) = \det \left(\sum_ {p \geq 0} \frac {1}{(p + 1) !} (\operatorname{ad} x) ^ {p}\right) \quad \text { for } x \in \mathfrak {h}.
\]

令 \(x, x_1, \ldots, x_m\) 为 \(\mathfrak{h}\) 中的元素。我们有

\[
(\exp^ {*} \omega_ {\mathrm{H}}) _ {x} (x _ {1}, \dots , x _ {m}) = (\omega_ {\mathrm{H}} (\exp x)) (\mathrm{T} _ {x} (\exp) (x _ {1}), \dots , \mathrm{T} _ {x} (\exp) (x _ {m})).
\]

将 \(\varpi(x) : \mathfrak{h} \to \mathfrak{h}\) 记为指数映射在 \(x\) 处的右微分（第 III 章，第 3 节，第 17 号，定义 8）；按定义，

\[
\mathrm{T} _ {x} (\exp) (y). (\exp x) ^ {- 1} = \varpi (x). y \quad \text { for   all } y \in \mathfrak {h}.
\]

由于形式 \(\omega_{H}\) 是右不变的，我们得到

\[
\begin{array}{r l} & (\omega_ {\mathrm{H}} (\exp x)) (\mathrm{T} _ {x} (\exp) (x _ {1}), \ldots , \mathrm{T} _ {x} (\exp) (x _ {m})) \\ & \quad = \omega_ {\mathrm{H}} (e) (\varpi (x). x _ {1}, \ldots , \varpi (x). x _ {m}) = (\det \varpi (x)) \omega_ {\mathfrak {h}} (x _ {1}, \ldots , x _ {m}); \end{array}
\]

因此，\(\exp^{*}\omega_{\mathrm{H}} = \lambda_{\mathfrak{h}}\omega_{\mathfrak{h}}\)，其中 \(\lambda_{\mathfrak{h}}(x) = \det \varpi(x) = \det \frac{\exp\operatorname{ad} x - 1}{\operatorname{ad} x}\)（第 III 章，第 6 节，第 4 号，命题 12）。

令 \(h \in H\)；由于 Ad h 是 h 的自同构，有

\[
\operatorname{ad} \left((\operatorname{Ad} h) (x)\right) = \operatorname{Ad} h \circ \operatorname{Ad} x \circ (\operatorname{Ad} h) ^ {- 1},
\]

所以 \(\lambda_{\mathfrak{h}}((\mathrm{Ad}h)(x)) = \lambda_{\mathfrak{h}}(x)\)。因此，函数 \(\lambda_{\mathfrak{h}}\) 在 \(\mathrm{Ad(H)}\) 下不变；命题得证。

注. 考虑与紧李群 G 关联的函数 \(\lambda_{g}\)；根据第 2 节第 1 号定理，要计算 \(\lambda_{g}\) 只需知道它在 t 上的值。但是，对于 \(x \in t\)，g 的自同态 ad x 是半单的，其特征值为 0（重数为 r），以及对于所有 \(\alpha \in \mathrm{R}(\mathrm{G}, \mathrm{T})\) 的 \(\delta(\alpha)(x)\)（重数为 1）。立即得到

\[
\lambda_ {\mathfrak {g}} (x) = \prod_ {\alpha \in \mathrm{R} (\mathrm{G}, \mathrm{T})} \frac {e ^ {\delta (\alpha) (x)} - 1}{\delta (\alpha) (x)} = \frac {\delta_ {\mathfrak {g}} (x)}{\pi_ {\mathfrak {g}} (x)}\tag{11}
\]

\[
\text { with } \delta_ {\mathfrak {g}} (x) = \delta_ {\mathrm{G}} (\exp x) \text { and } \pi_ {\mathfrak {g}} (x) = \prod_ {\alpha \in \mathrm{R} (\mathrm{G}, \mathrm{T})} \delta (\alpha) (x) = \det \operatorname{ad} _ {\mathfrak {g} / \mathfrak {t}} (x).
\]

令 \(\omega_{\mathrm{G / T}}\) 为次数 \(n - r\) 的 \(\mathrm{G / T}\) 上不变微分形式，令 \(\omega_{\mathfrak{t}}\) 为次数 \(r\) 的 \(\mathfrak{t}\) 上平移不变微分形式。按第 1 号的记号，将 \(\omega_{\mathrm{G / T}} \cap \omega_{\mathfrak{t}}\) 记为唯一的平移不变微分形式 \(\omega_{\mathfrak{g}}\)，其次数为 \(n\)，定义在 \(\mathfrak{g}\) 上，且满足 \(\omega_{\mathfrak{g}}(0) = \omega_{\mathrm{G / T}}(\bar{e}) \cap \omega_{\mathfrak{t}}(0)\)。

最后，将 \(\psi:(\mathrm{G}/\mathrm{T})\times\mathfrak{t}\to\mathfrak{g}\) 记为由映射 \((g,x)\mapsto(\operatorname{Ad}g)(x)\) 从 \(G\times t\) 到 g 通过取商诱导的流形态射。

命题 3. 设 \(\omega_{\mathfrak{g}}\)、\(\omega_{\mathfrak{t}}\)、\(\omega_{\mathrm{G / T}}\) 分别是 \(\mathfrak{g}\)、\(\mathfrak{t}\)、\(\mathrm{G / T}\) 上次数分别为 \(n, r, n - r\) 的不变微分形式。若 \(\omega_{\mathfrak{g}} = \omega_{\mathrm{G / T}} \cap \omega_{\mathfrak{t}}\)，则有

\[
\psi^ {*} \omega_ {\mathfrak {g}} = \omega_ {\mathrm{G/T}} \wedge \pi_ {\mathfrak {g}} \omega_ {\mathfrak {t}}
\]

其中 \(\pi_{\mathfrak{g}}\) 是 \(\mathfrak{t}\) 上由 \(\pi_{\mathfrak{g}}(x) = \prod_{\alpha \in \mathrm{R}(\mathrm{G},\mathrm{T})}\delta (\alpha)(x)\) 定义的函数。

将 \(\omega_{G}\)（分别为 \(\omega_{T}\)）记为 G（分别为 T）上最高次数的不变微分形式，它们在原点处分别与 \(\omega_{g}\)（分别为 \(\omega_{t}\)）相同。考虑交换图

\[
\begin{array}{c c c} (\mathrm{G} / \mathrm{T}) \times \mathfrak {t} & \stackrel {{\psi}} {{\longrightarrow}} & \mathfrak {g} \\ \Big \downarrow^ {(\mathrm{Id}, \exp_ {\mathrm{T}})} & & \Big \downarrow^ {\exp_ {\mathrm{G}}} . \\ (\mathrm{G} / \mathrm{T}) \times \mathrm{T} & \stackrel {{f}} {{\longrightarrow}} & \mathrm{G} \end{array}
\]

根据第 2 号的命题 1 以及关系 \(\exp_{T}^{*}\omega_{T} = \omega_{t}\)，得到等式

\[
\psi^ {*} \mathrm{exp} _ {\mathrm{G}} ^ {*} \omega_ {\mathrm{G}} = \omega_ {\mathrm{G/T}} \wedge \delta_ {\mathfrak {g}} \omega_ {\mathfrak {t}}.
\]

根据命题 2，\(\psi^{*}\exp_{\mathrm{G}}^{*}\omega_{\mathrm{G}} = (\psi^{*}\lambda_{\mathfrak{g}})\psi^{*}\omega_{\mathfrak{g}}\)。由于函数 \(\lambda_{\mathfrak{g}}\) 在 \(\operatorname{Ad}(\mathrm{G})\) 下不变，有

\[
(\psi^ {*} \lambda_ {\mathfrak {g}}) (\bar {g}, x) = \lambda_ {\mathfrak {g}} (x) = \frac {\delta_ {\mathfrak {g}} (x)}{\pi_ {\mathfrak {g}} (x)} \quad \text { for } \bar {g} \in \mathrm{G/T}, x \in \mathfrak {t}.
\]

由此可知，形式 \(\psi^{*}\omega_{\mathrm{G}}(\bar{g},x)\) 与 \(\omega_{\mathrm{G / T}}(\bar{g})\wedge \pi_{\mathfrak{g}}(x)\omega_{\mathfrak{t}}(x)\) 在 \(\delta_{\mathfrak{g}}(x)\) 非零处相同，也就是在稠密开子集 \((\mathrm{G / T})\times \mathfrak{t}_r\) 上相同；因此它们处处相等，命题得证。

在 G 和 T 上分别选取最高次数的不变微分形式 \(\omega_{G}\) 和 \(\omega_{T}\)，使得 \(|\omega_{G}| = dg\) 且 \(|\omega_{T}| = dt\)；将 \(\omega_{g}\)（分别为 \(\omega_{t}\)）记为 g（分别为 t）上在原点处与 \(\omega_{\mathrm{G}}(e)\)（分别为 \(\omega_{\mathrm{T}}(e)\)）相同的平移不变微分形式，并将 Haar 测度 \(|\omega_{g}|\)（分别为 \(|\omega_{t}|\)）记为 dz（分别为 dx）。作与第 2 号类似并作相应修改的推理，得到以下命题：

命题 4. 测度 \(dz\)（定义在 \(\mathfrak{g}\) 上）是映射 \((g,x)\mapsto (\operatorname {Ad}g)(x)\) 从 \(\mathrm{G}\times \mathfrak{t}\) 到 \(\mathfrak{g}\) 的真映射下测度 \(dg\otimes \frac{1}{w(\mathrm{G})}\pi_{\mathfrak{g}}dx\) 的像。

第 2 号的推论 1 至 3 以及备注 1 至 3 的类似命题的叙述和证明留给读者。例如，设 \(\varphi\) 是 \(\mathfrak{g}\) 上的可积函数（取值于 Banach 空间或 \(\overline{\mathbf{R}}\)）；则

\[
\int_ {\mathfrak {g}} \varphi (z) d z = \frac {1}{w (\mathrm{G})} \int_ {\mathfrak {t}} \left(\int_ {\mathrm{G}} \varphi ((\mathrm{Ad} g) x) d g\right) \pi_ {\mathfrak {g}} (x) d x,\tag{12}
\]

特别地，若 \(\varphi\) 在 \(\operatorname{Ad}(\mathbf{G})\) 下不变，则

\[
\int_ {\mathfrak {g}} \varphi (z) d z = \frac {1}{w (\mathrm{G})} \int_ {\mathfrak {t}} \varphi (x) \pi_ {\mathfrak {g}} (x) d x.\tag{13}
\]

\subsection{向量丛截面的积分}

在本号及下一号中，用 X 表示一个 \(C^{r}\) 类的实流形 \((1 \leq r \leq \infty)\)，且它局部有限维。

令 Y 为一个 \(C^{r}\) 类流形。若 \(r < \infty\)，考虑映射 \(f \mapsto j^{r}(f)\)，它从 \(\mathcal{C}^{r}(\mathrm{X};\mathrm{Y})\) 到 \(\mathcal{C}(\mathrm{X};\mathrm{J}^{r}(\mathrm{X},\mathrm{Y}))\)（《微分与解析流形，结果》，12.3.7）。\(\mathcal{C}(\mathrm{X};\mathrm{J}^{r}(\mathrm{X},\mathrm{Y}))\) 上的紧收敛拓扑在该映射下的逆像称为紧 \(C^{r}\) 收敛拓扑在 \(\mathcal{C}^{r}(\mathrm{X};\mathrm{Y})\) 上的拓扑；它是关于 K 的一致 \(C^{r}\) 收敛拓扑的上确界（《微分与解析流形，结果》，12.3.10），其中 K 遍历 X 的紧子集。

当 \(r = \infty\) 时，将紧 \(C^{\infty}\) 收敛拓扑定义在 \(\mathcal{C}^{\infty}(X;Y)\) 上，使其成为紧 \(C^{k}\) 收敛拓扑的上确界；换言之，它是使典则嵌入 \(\mathcal{C}^{\infty}(X;Y) \to \mathcal{C}^{k}(X;Y)\) 对所有 \(0 \leq k < \infty\) 均连续的最粗拓扑。

令 \(\mathrm{E}\) 为以 \(\mathrm{X}\) 为底空间的 \(\mathrm{C}^r\) 类实向量丛，令 \(\mathcal{S}^r(\mathrm{X};\mathrm{E})\) 为 \(\mathrm{E}\) 的 \(\mathrm{C}^r\) 类截面的向量空间。在本号中，赋予 \(\mathcal{S}^r(\mathrm{X};\mathrm{E})\) 由紧 \(\mathrm{C}^r\) 收敛拓扑在 \(\mathcal{C}^r(\mathrm{X};\mathrm{E})\) 上诱导的拓扑，也称为紧 \(\mathrm{C}^r\) 收敛拓扑；它使 \(\mathcal{S}^r (\mathrm{X};\mathrm{E})\) 成为完备、分离的局部凸拓扑向量空间（参见《微分与解析流形，结果》，15.3.1 以及《谱理论》（编写中））。

现在令 H 为李群，\(m: H \times X \to X\) 为 \(C^{r}\) 类的左作用律；置 \(hx = m(h, x)\)，其中 \(h \in H\)、\(x \in X\)。令 E 为以 X 为底空间的 \(C^{r}\) 类向量 H-丛（第 III 章，第 1 节，第 8 号，定义 4）。对于 \(s \in \mathcal{S}^{r}(X; E)\) 和 \(h \in H\)，将 \(^{h}s\) 记为 E 的截面 \(x \mapsto h.s(h^{-1}x)\)；映射 \((h, s) \mapsto ^{h}s\) 是 H 在空间 \(\mathcal{S}^{r}(X; E)\) 上的作用律。

引理 4. 作用律 \(\mathrm{H} \times \mathcal{S}^r (\mathrm{X};\mathrm{E}) \to \mathcal{S}^r (\mathrm{X};\mathrm{E})\) 是连续的。

根据 \(\mathcal{S}^{r}(\mathrm{X};\mathrm{E})\) 的拓扑定义以及《一般拓扑学》第 X 章第 3 节第 4 号定理 3，只需证明：对于任意整数 \(k \leq r\)，映射 \(f : H \times X \times \mathcal{S}^{k}(\mathrm{X};\mathrm{E}) \to \mathrm{J}^{k}(\mathrm{X};\mathrm{E})\) 满足 \(f(h,x,s) = j_{x}^{k}(^{h}s)\) 且是连续的。对于 \(h \in H\)，将 \(\tau_{h}\)（分别为 \(\theta_{h}\)）记为 X（分别为 E）上的自同构 \(x \mapsto hx\)。定义映射

\[
f _ {1}: \mathrm{H} \times \mathrm{X} \to \mathrm{J} ^ {k} (\mathrm{X}, \mathrm{X})
\]

\[
f _ {2}: \mathrm{H} \times \mathrm{E} \rightarrow \mathrm{J} ^ {k} (\mathrm{E}, \mathrm{E})
\]

\[
g: \mathrm{H} \times \mathrm{X} \times \mathcal {S} ^ {k} (\mathrm{X}; \mathrm{E}) \to \mathrm{J} ^ {k} (\mathrm{X}, \mathrm{E})
\]

由 \(f_{1}(h,x) = j_{x}^{k}(\tau_{h}),f_{2}(h,v) = j_{v}^{k}(\theta_{h}),g(h,x,s) = j_{hx}^{k}(s)\) 定义。我们有

\[
f (h, x, s) = f _ {2} (h, s (h ^ {- 1} x)) \circ g (h ^ {- 1}, x, s) \circ f _ {1} (h ^ {- 1}, x),
\]

因此，根据《微分与解析流形，结果》12.3.6，只需证明 \(f_{1}, f_{2}\) 和 \(g\) 连续。

现在 \(g\) 是复合映射

\[
\begin{array}{c c c} \mathrm{H} \times \mathrm{X} \times \mathcal {S} ^ {k} (\mathrm{X}; \mathrm{E}) & \xrightarrow {(m , \mathrm{Id})} & \mathrm{X} \times \mathcal {S} ^ {k} (\mathrm{X}; \mathrm{E}) \\ & \xrightarrow {(\mathrm{Id} , j ^ {k})} & \mathrm{X} \times \mathcal {C} (\mathrm{X}; \mathrm{J} ^ {k} (\mathrm{X}, \mathrm{E})) \quad \xrightarrow {\varepsilon} \quad \mathrm{J} ^ {k} (\mathrm{X}, \mathrm{E}) \end{array}
\]

其中 \(\varepsilon(x, u) = u(x)\)；由于映射 \(\varepsilon\) 连续（《一般拓扑学》第 X 章第 3 节第 4 号定理 3 的推论 1），\(g\) 连续。

令 \((h_{0}, x_{0}) \in \mathrm{H} \times \mathrm{X}\)；我们证明 \(f_{1}\) 在 \((h_{0}, x_{0})\) 处连续。存在图表 \((\mathrm{U}, \psi, \mathrm{F})\) 和 \((\mathrm{V}, \chi, \mathrm{F}')\)（均为 \(\mathrm{X}\) 的图表），以及开子集 \(\Omega\)（属于 \(\mathrm{H}\)），使得 \(x_{0} \in \mathrm{U}, h_{0} \in \Omega\) 且 \(m(\Omega \times \mathrm{U}) \subset \mathrm{V}\)。利用这些图表中的 \(\mathrm{J}^{k}(\mathrm{X}, \mathrm{X})\) 表达式，问题归结为证明：对于 \(1 \leq l \leq k\)，在 \((h_{0}, x_{0})\) 处映射 \((h, x) \mapsto \Delta_{x}^{l}(\tau_{h})\) 从 \(\Omega \times \mathrm{U}\) 到 \(\mathrm{P}_{l}(\mathrm{F}; \mathrm{F}')\) 连续，其中 \(\Delta_{x}^{l}(\tau_{h})(v) = \frac{1}{l!}\mathrm{D}^{l}\tau_{h}(x).v\) 对 \(v \in \mathrm{F}\) 成立（《微分与解析流形，结果》，12.2）。但是，\(\mathrm{D}^{l}\tau_{h}(x)\) 正是 \(l\) 阶偏导数 \(m(h, x)\) 关于 \(x\) 的，按假设它是连续的；因此 \(f_{1}\) 连续。\(f_{2}\) 连续性的证明类似，故引理得证。

命题 5. 假设群 H 是紧的，并将 H 上总质量为 1 的 Haar 测度记为 dh。令 s 为 E 的 \(C^{r}\) 类截面。

将 \(s^{\sharp}\) 记为向量积分 \(\int_{\mathrm{H}}{}^{h}s\,dh\)。则 \(s^{\sharp}\) 是 \(\mathrm{E}\) 的 \(\mathrm{C}^r\) 类截面，在 \(\mathrm{H}\) 下不变；对于 \(x \in \mathrm{X}\)，有 \(s^{\sharp}(x) = \int_{\mathrm{H}}hs(h^{-1}x)\,dh \in \mathrm{E}_x\)。将 \(s \mapsto s^{\sharp}\) 记为 \(\mathcal{S}^r(\mathrm{X};\mathrm{E})\) 上的自同态，它是到 H-不变截面子空间上的投影。

考虑映射 \(h \mapsto {}^{h}s\) 到 \(\mathcal{S}^{r}(\mathrm{X};\mathrm{E})\)；由引理 4 它是连续的。由于空间 \(\mathcal{S}^{r}(\mathrm{X};\mathrm{E})\) 分离且完备，积分 \(s^{\sharp} = \int_{H}{}^{h}s\,dh\) 属于 \(\mathcal{S}^{r}(\mathrm{X};\mathrm{E})\)（积分，第 III 章，第 3 节，第 3 号，推论 2）。线性映射 \(s \mapsto s(x)\) 从 \(\mathcal{S}^{r}(\mathrm{X};\mathrm{E})\) 到 \(E_{x}\) 且连续，因此 \(s^{\sharp}(x) = \int_{H}{}^{h}s(x)\,dh\) 对所有 \(x \in X\) 成立。显然 \(s^{\sharp}\) 在 H 下不变；若 s 是 H-不变截面，则 \(s^{\sharp} = s\)，故最后一个断言成立。

推论 1. 令 \(\mathrm{F}\) 为 Banach 空间，\(\rho : \mathrm{H} \to \mathbf{GL}(\mathrm{F})\) 为解析线性表示，\(f \in \mathcal{C}^r(\mathrm{X};\mathrm{F})\)。对于 \(x \in \mathrm{X}\)，置

\[
f ^ {\sharp} (x) = \int_ {\mathrm{H}} \rho (h). f (h ^ {- 1} x) d h.
\]

则 \(f^{\sharp}\) 是 \(\mathbf{C}^r\) 类的从 \(\mathbf{X}\) 到 \(\mathbf{F}\) 的态射，并且与 \(\mathbf{H}\) 的作用相容；对于 \(x \in \mathbf{X}\)，有（其中 \(\tau_h\) 表示自同构 \(x \mapsto hx\) 的 \(\mathbf{X}\)）

\[
d _ {x} f ^ {\sharp} = \int_ {\mathrm{H}} (\rho (h) \circ d _ {h ^ {- 1} x} f \circ \mathrm{T} _ {x} (\tau_ {h ^ {- 1}})) d h \in \mathscr {L} (\mathrm{T} _ {x} (\mathrm{X}); \mathrm{F}).\tag{14}
\]

第一断言由将命题应用于丛 \(X \times F\) 得到，其中赋予该丛以作用律 \((h; (x, f)) \mapsto (hx, \rho(h).f)\)。第二断言由积分，第 III 章第 3 节第 2 号命题 2 得到：对向量积分 \(f^{\sharp}\) 应用同态 \(d_{x} : \mathcal{C}^{r}(\mathrm{X};\mathrm{F}) \to \mathcal{L}(\mathrm{T}_{x}(\mathrm{X});\mathrm{F})\)，而该同态根据紧 \(C^{r}\) 收敛拓扑的定义是连续的。

推论 2. 令 F 为 Banach 空间，\(f \in \mathcal{C}^{r}(X;F)\)；置

\[
f ^ {\sharp} (x) = \int_ {\mathrm{H}} f (h x) d h
\]

对于 \(x \in \mathrm{X}\)。函数 \(f^{\sharp}\) 是 \(\mathbf{C}^r\) 类的，并且 \(f^{\sharp}(hx) = f^{\sharp}(x)\) 对于 \(x \in \mathrm{X}\)、\(h \in \mathrm{H}\) 成立。

推论 3. 令 F 为 Banach 空间，p 为满足 \(\geq 0\) 的整数，令 \(^{k}\Omega^{p}(\mathrm{X};\mathrm{F})\) 为 X 上取值于 F 且属于 \(C^{k}\) 类的 p 次微分形式空间（ \(2 \leq k + 1 \leq r\)）。对于 \(\omega \in {}^{k}\Omega^{p}(\mathrm{X};\mathrm{F})\)，置 \(\omega^{\sharp} = \int_{H} \tau_{h}^{*}\omega dh\)。则映射 \(\omega \mapsto \omega^{\sharp}\) 是 \(^{k}\Omega^{p}(\mathrm{X};\mathrm{F})\) 上的投影，其像为 H-不变形式子空间；有 \(d(\omega^{\sharp}) = (d\omega)^{\sharp}\) 对所有 \(\omega \in {}^{k}\Omega^{p}(\mathrm{X};\mathrm{F})\) 成立。

第一断言由将命题应用于向量 H-丛 \(\operatorname{Alt}^{p}(\mathrm{T}(\mathrm{X});\mathrm{F})\) 得到（第 III 章，第 1 节，第 8 号，例）。要证明第二断言，根据积分，第 III 章第 3 节第 2 号命题 2，只需证明映射 \(d:{}^{k}\varOmega^{p}(\mathrm{X};\mathrm{F})\to{}^{k-1}\varOmega^{p+1}(\mathrm{X};\mathrm{F})\) 是连续的，其中第一（分别为第二）空间赋予紧 \(C^{k}\) 收敛拓扑（分别为 \(C^{k-1}\) 收敛拓扑）。但这立即由利用半范数给出的这些拓扑的定义（《谱理论》（编写中））以及 d 是阶数 \(\leq 1\) 的微分算子这一事实得到（《微分与解析流形，结果》，14.4.2）。

\subsection{不变微分形式}

令 X 为局部有限维的 \(C^{\infty}\) 类实流形，并令 \((g, x) \mapsto gx\) 为连通紧李群 G 在 X 上的 \(C^{\infty}\) 类左作用律。对于 \(g \in G\)，将 \(\tau_{g}\) 记为 X 的自同构 \(x \mapsto gx\)。将 \(\Omega(X)\) 记为 X 上 \(C^{\infty}\) 类实微分形式的代数（《微分与解析流形，结果》，8.3.1）。

对于任意 \(\xi\) 属于 \(\mathfrak{g}\)，将相应的 X 上向量场记为 \(D_{\xi}\)（第 III 章，第 3 节，第 5 号），并将相应的算子 \(\theta(\xi), i(\xi)\) 记为 \(\Omega(X)\) 上的算子，于是有公式（《微分与解析流形，结果》，8.4.5 和 8.4.7）

\[
\theta (\xi) \omega = d (i (\xi) \omega) + i (\xi) d \omega\tag{15}
\]

\[
\frac {d}{d t} (\tau_ {\mathrm{exp} t \xi} ^ {*} \omega) = \tau_ {\mathrm{exp} t \xi} ^ {*} (\theta (\xi) \omega).\tag{16}
\]

若微分形式 \(\omega \in \Omega(\mathrm{X})\) 满足 \(\tau_g^*\omega = \omega\) 对所有 \(g \in \mathrm{G}\) 成立，则称其为不变的；由公式（16），这等价于 \(\theta(\xi)\omega = 0\) 对所有 \(\xi \in \mathfrak{g}\) 成立。将 X 上不变微分形式的空间记为 \(\Omega(\mathrm{X})^{\mathrm{G}}\)；若 \(\omega \in \Omega(\mathrm{X})^{\mathrm{G}}\)，则有 \(d\omega \in \Omega(\mathrm{X})^{\mathrm{G}}\)，所以 \(\Omega(\mathrm{X})^{\mathrm{G}}\) 是复形 \((\Omega(\mathrm{X}), d)\) 的子复形。

定理 2. 典则嵌入 \(\iota : \Omega(\mathrm{X})^{\mathrm{G}} \to \Omega(\mathrm{X})\) 是复形的同伦等价（《代数》，第 X 章，第 33 页，定义 5）；映射 \(\omega \mapsto \omega^{\sharp} = \int_{\mathrm{G}} \tau_g^* \omega dg\) 是一个同伦等价，在同伦意义下为其逆。特别地，映射 \(\mathrm{H}(\iota) : \mathrm{H}(\Omega(\mathrm{X})^{\mathrm{G}}) \to \mathrm{H}(\Omega(\mathrm{X}))\) 是双射。

由第 4 号的推论 3，映射 \(\omega \mapsto \omega^{\sharp}\) 是从 \(\Omega(\mathrm{X})\) 到 \(\Omega(\mathrm{X})^{\mathrm{G}}\) 的复形态射，并且在子复形 \(\Omega(\mathrm{X})^{\mathrm{G}}\) 上诱导恒等映射；因此，要证明定理，只需构造一个分次同态 \(s: \Omega(\mathrm{X}) \to \Omega(\mathrm{X})\)，其次数为 \(-1\)，使得

\[
\omega^ {\sharp} - \omega = (d \circ s + s \circ d) (\omega) \quad \text { for   all } \omega \in \Omega (\mathrm{X}).\tag{17}
\]

由《积分》第 IX 章第 2 节第 4 号的引理 1 以及第 2 节第 2 号的备注 1，存在正测度 \(d\xi\)，定义在具有紧支集的 \(\mathfrak{g}\) 上，其在指数映射下的像等于 \(dg\)。对于 \(\omega \in \Omega(X)\)，置

\[
s (\omega) = \int_ {0} ^ {1} \left\{\int_ {\mathfrak {g}} \tau_ {\exp t \xi} ^ {*} (i (\xi) \omega). d \xi \right\} d t;
\]

我们需证明公式（17）成立。如同第 4 号推论 1 的证明，我们验证公式

\[
d s (\omega) = \int_ {0} ^ {1} \left\{\int_ {\mathfrak {g}} \tau_ {\exp t \xi} ^ {*} d (i (\xi) \omega). d \xi \right\} d t.
\]

现在由公式（15）和（16）得到等式

\[
\begin{array}{r l} & {d s (\omega) + s (d \omega) = \int_ {0} ^ {1} \left\{\int_ {\mathfrak {g}} \tau_ {\exp t \xi} ^ {*} (d (i (\xi) \omega) + i (\xi) d \omega). d \xi \right\} d t} \\ & {\qquad = \int_ {0} ^ {1} \left\{\int_ {\mathfrak {g}} \tau_ {\exp t \xi} ^ {*} (\theta (\xi) \omega). d \xi \right\} d t} \\ & {\qquad = \int_ {\mathfrak {g}} \left\{\int_ {0} ^ {1} \frac {d}{d t} (\tau_ {\exp t \xi} ^ {*} \omega) d t \right\} d \xi} \\ & {\qquad = \int_ {\mathfrak {g}} (\tau_ {\exp \xi} ^ {*} \omega - \omega) d \xi} \\ & {\qquad = \omega^ {\sharp} - \omega ,} \end{array}
\]

定理 2 得证。

将定理应用于 X = G 且 G 通过左平移作用的情形。回忆（第 III 章，第 3 节，第 13 号，命题 50），将 G 上的微分形式与其在单位元处的值对应，给出从 \(\Omega(\mathrm{G})^{\mathrm{G}}\) 到交错多重线性形式的分次代数 \(\operatorname{Alt}(\mathfrak{g})\) 的同构。利用此同构将 \(\Omega(\mathrm{G})^{\mathrm{G}}\) 与 \(\operatorname{Alt}(\mathfrak{g})\) 识别。于是算子 d 由下式给出（第 III 章，第 3 节，第 14 号，命题 51）

\[
\begin{array}{l} d \omega (a _ {1}, \ldots , a _ {p + 1}) \\ = \sum_ {i <   j} (- 1) ^ {i + j} \omega ([ a _ {i}, a _ {j} ], a _ {1}, \ldots , a _ {i - 1}, a _ {i + 1}, \ldots , a _ {j - 1}, a _ {j + 1}, \ldots , a _ {p + 1}) \end{array}
\]

其中 \(\omega\) 属于 \(\mathrm{Alt}^p (\mathfrak{g})\)，且 \(a_1,\ldots ,a_{p + 1}\) 属于 \(\mathfrak{g}\)。

对于 \(\xi \in \mathfrak{g}\)，令相应的左不变向量场为 \(\mathrm{L}_{\xi}\)（利用 G 通过右平移在自身上的作用定义，参见第 III 章，第 3 节，第 6 号）。算子 \(\theta (\mathrm{L}_{\xi}), i(\mathrm{L}_{\xi})\) 与左平移所定义的 G 在 \(\Omega (\mathrm{G})\) 上的作用交换，因而诱导算子 \(\theta (\xi), i(\xi)\)，它们作用在 \(\Omega (\mathrm{G})^{\mathrm{G}}\) 上；在上述识别下，它们由下列公式表示（《微分与解析流形，结果》，8.3.2 和 8.4.2）

\[
(\theta (\xi) \omega) (a _ {1}, \dots , a _ {p}) = - \sum_ {i} \omega (a _ {1}, \dots , a _ {i - 1}, [ \xi , a _ {i} ], a _ {i + 1}, \dots , a _ {p})
\]

\[
(i (\xi) \omega) (a _ {1}, \ldots , a _ {p - 1}) = \omega (\xi , a _ {1}, \ldots , a _ {p - 1})
\]

其中 \(\omega\) 属于 \(\operatorname{Alt}^p (\mathfrak{g})\)，且 \(a_1,\ldots ,a_p\) 属于 \(\mathfrak{g}\)。

由双不变形式组成的子复形 \({}^{\mathrm{G}}\varOmega(\mathrm{G})^{\mathrm{G}}\)（第 III 章，第 3 节，第 13 号）可与子复形 \(\operatorname{Alt}(\mathfrak{g})^{\mathrm{G}}\) 识别，后者由在伴随表示下不变的 g 上交错多重线性形式组成（即，满足 \(\theta(\xi)\omega=0\) 对所有 \(\xi\in\mathfrak{g}\) 成立）。因此有如下复形交换图

\[
\begin{array}{c c c c c} ^ {\mathrm{G}} \varOmega (\mathrm{G}) ^ {\mathrm{G}} & \longrightarrow & \varOmega (\mathrm{G}) ^ {\mathrm{G}} & \longrightarrow & \varOmega (\mathrm{G}) \\ \Big \downarrow & & \Big \downarrow & & \\ \operatorname{Alt} (\mathfrak {g}) ^ {\mathrm{G}} & \longrightarrow & \operatorname{Alt} (\mathfrak {g}) & & \end{array}\tag{18}
\]

其中水平箭头是典则嵌入，竖直箭头是由映射 \(\omega\mapsto\omega(e)\) 诱导的同构。

推论 1. a) 在图（18）中，所有态射都是同伦等价。b) 令 \(\omega \in \operatorname{Alt}(\mathfrak{g})\)。则 \(\omega\) 属于 \(\operatorname{Alt}(\mathfrak{g})^{\mathrm{G}}\) 当且仅当 \(d\omega = 0\) 且 \(d(i(\xi)\omega) = 0\) 对所有 \(\xi \in g\) 成立。复形 \(\operatorname{Alt}(\mathfrak{g})^{\mathrm{G}}\) 的微分为零。c) 分次向量空间 \(\mathrm{H}(\Omega(\mathrm{G}))\) 同构于 \(\operatorname{Alt}(\mathfrak{g})^{\mathrm{G}}\)。

将定理应用于 G 通过左平移在 G 上的作用（分别应用于 G × G 通过 \(((g,h);x)\mapsto gxh^{-1}\) 在 G 上的作用），得到典则嵌入 \(\Omega(\mathrm{G})^{\mathrm{G}}\to\Omega(\mathrm{G})\)（分别为 \({}^{G}\Omega(\mathrm{G})^{\mathrm{G}}\to\Omega(\mathrm{G})\)）是同伦等价；根据《代数》第 X 章第 34 页的推论，断言 a) 成立。

我们证明 \(b)\)。根据第 III 章第 3 节第 14 号命题 51，有 \(d\alpha = -d\alpha\)，即 \(d\alpha = 0\)，对于微分形式 \(\alpha\)（定义在 \(G\) 上且左、右不变）成立。因此，若 \(\omega \in \mathrm{Alt}(\mathfrak{g})^G\)，则 \(d\omega = 0\)，从而 \(d(i(\xi)\omega) = \theta(\xi)\omega - i(\xi)d\omega = 0\)。反之，若 \(d\omega = 0\) 且 \(d(i(\xi)d\omega) = 0\)，则 \(\theta(\xi)\omega = 0\)。断言 \(c)\) 由 \(a)\) 和 \(b)\) 得出。

注. 考虑子复形 \(\mathrm{Z}(\varOmega(\mathrm{G}))\) 和 \(\mathrm{B}(\varOmega(\mathrm{G}))\)（它们属于 \(\varOmega(\mathrm{G})\)）（《代数》，第 X 章，第 25 页）。由两个形式乘积的微分公式（《微分与解析流形，结果》，8.3.5）可知，\(\mathrm{Z}(\varOmega(\mathrm{G}))\) 是 \(\varOmega(\mathrm{G})\) 的子代数，而 \(\mathrm{B}(\varOmega(\mathrm{G}))\) 是 \(\mathrm{Z}(\varOmega(\mathrm{G}))\) 的理想；因此，外积在 \(\mathrm{H}(\varOmega(\mathrm{G}))\) 上诱导分次代数结构。上述结果给出分次代数的同构 \(\mathrm{H}(\varOmega(\mathrm{G})) \to \mathrm{Alt}(\mathfrak{g})^{\mathrm{G}}\)。

令 H 为 G 的闭子群；将定理 2 应用于 X = G/H。根据第 III 章第 1 节第 8 号命题 17 的推论 1，G/H 上的 G-不变微分形式可与 \(\operatorname{Alt}(\mathrm{T}_{e}(\mathrm{G}/\mathrm{H}))\) 的 H-不变元素识别，即与 \(\operatorname{Alt}(\mathfrak{g})\) 中在 H 下不变且被算子 \(i(\xi)\) 对所有 \(\xi \in \mathrm{L}(\mathrm{H})\) 湮没的元素识别。因此：

推论 2. 设 \(\mathrm{H}\) 为 \(\mathrm{G}\) 的闭子群。

\begin{enumerate}
\def\labelenumi{\alph{enumi})}
\item
  典则嵌入 \(\Omega (\mathrm{G / H})^{\mathrm{G}}\to \Omega (\mathrm{G / H})\) 是同伦等价。
\item
  复形 \(\Omega(\mathrm{G}/\mathrm{H})^{\mathrm{G}}\) 可与 \(\operatorname{Alt}(\mathfrak{g})\) 的如下子复形识别：它由元素 \(\omega\)（属于 \(\operatorname{Alt}(\mathfrak{g})\)）在 H 的伴随表示下不变，且满足 \(i(\xi)\omega = 0\) 对所有 \(\xi \in \operatorname{L}(\operatorname{H})\) 成立的元素组成。如果 H 还是连通的，则该子复形由 \(\omega \in \operatorname{Alt}(\mathfrak{g})\) 组成，且满足 \(\theta(\xi)\omega = 0\) 和 \(i(\xi)\omega = 0\) 对所有 \(\xi \in \operatorname{L}(\operatorname{H})\) 成立。
\end{enumerate}

\section{连通紧李群的不可约表示}

沿用第 6 节的记号。G 的一个表示是从 G 到群 \(\mathbf{GL}(\mathrm{V})\) 的连续（因而解析）同态，其中 V 是有限维复向量空间。G 的每个表示都是半单的（第 1 节，第 1 号）。

选取 \(\mathfrak{t}\) 的一个房 C（ \(\S 5\)，第 2 号），并置 \(\Gamma(\mathrm{T})_{++} = \overline{\mathrm{C}} \cap \Gamma(\mathrm{T})\)。

\subsection{支配特征标}

将 \(X_{+}\) 定义为如下集合：其元素为 \(\lambda\)，且属于 \(\mathrm{X(T)}\)，并满足 \(\langle\lambda,x\rangle\geq0\) 对所有 \(x\in\Gamma(\mathrm{T})_{++}\) 成立；也就是线性形式 \(\delta(\lambda):\mathfrak{t}_{\mathbf{C}}\to\mathbf{C}\) 将 t 的房 C 映到 \(iR_{+}\)。

赋予 X(T) 以序群结构，其中正元素是 \(X_{+}\) 中的元素；置 \(R_{+} = R(G, T) \cap X_{+}\)，\(R_{-} = -R_{+}\)。\(R_{+}\) 的元素称为正根，\(R_{-}\) 的元素称为负根；每个根或者是正根，或者是负根（第 VI 章，第 1 节，第 6 号，定理 3）。不可能写成两个正根之和的正根称为单根；每个正根都是单根之和（同上）；单根构成 X(T) 中由根生成的子群的一个基，该子群可与 X(T/C(G)) 识别（第 4 节，第 4 号）；关于单根的反射生成 Weyl 群 \(W = W_{G}(T)\)（第 VI 章，第 1 节，第 5 号，定理 2）。

引理 1. 令 \(\lambda\) 为 \(\mathrm{X(T)}\) 的一个元素。以下条件等价：

\begin{enumerate}
\def\labelenumi{(\roman{enumi})}
\item
  \(\lambda - w(\lambda) \geq 0\)（分别为 \(>0\)）对所有 \(w \in \mathrm{W}\) 满足 \(w \neq 1\) 成立；
\item
  当 \(w \in \mathrm{W}\) 且 \(w \neq 1\) 时，\(\lambda - w(\lambda)\) 是单根的具有正系数的线性组合（分别为不全为零的正系数线性组合）；
\item
  \(\langle \lambda, K_{\alpha} \rangle \geq 0\)（分别为 \(>0\)）对每个正根 \(\alpha\) 成立；
\item
  \(\langle\lambda,K_{\alpha}\rangle\geq0\)（分别为 >0）对每个单根 \(\alpha\) 成立。
\end{enumerate}

条件（iii）与（iv）的等价性是立即的。由于 \(K_{\alpha}\) 的集合可与 \(\mathrm{R}(\mathrm{G},\mathrm{T})\) 的逆根系识别（第 4 节，第 5 号），（i）与（iii）的等价性由第 VI 章第 1 节第 6 号命题 18 及其推论得到。(ii) \(\Rightarrow\) (i) 是平凡的，反向蕴含由同上引文得到。

将 \(X_{++}\) 记为 \(\mathrm{X(T)}\) 中满足 \(\langle\lambda,K_{\alpha}\rangle\geq0\) 对每个正根 \(\alpha\) 成立的元素集合。\(X_{++}\) 的元素称为支配的。它们构成 W 在 \(\mathrm{X(T)}\) 上作用的基本域（第 VI 章，第 1 节，第 10 号）。有 \(X_{++}\subset X_{+}\)。

若 G 单连通，则对于每个单根 \(\alpha\)，存在 X(T) 的元素 \(\varpi_{\alpha}\)，使得 \(\langle\varpi_{\beta}, K_{\alpha}\rangle = \delta_{\alpha\beta}\) 对每个单根 \(\beta\) 成立，即 \(s_{\alpha}(\varpi_{\alpha}) = \varpi_{\alpha} - \alpha\)，且 \(s_{\beta}(\varpi_{\alpha}) = \varpi_{\alpha}\) 对每个单根 \(\beta \neq \alpha\) 成立；\(\varpi_{\alpha}\) 称为基本支配权；它们构成交换群 \(\mathrm{X(T)}\) 和交换幺半群 \(\mathrm{X}_{++}\) 的基；更精确地，每个元素 \(\lambda\) 属于 \(\mathrm{X(T)}\)，都可写成 \(\lambda = \sum_{\alpha} \langle \lambda, K_{\alpha} \rangle \varpi_{\alpha}\) 的形式。

将 \(\rho\) 记为 \(\mathrm{X(T)}\otimes \mathbf{Q}\) 中满足下式的元素：

\[
2 \rho = \sum_ {\alpha \in \mathrm{R} _ {+}} \alpha .
\]

\(\langle\rho,K_{\alpha}\rangle=1\) 对每个单根 \(\alpha\) 成立（第 VI 章，第 1 节，第 10 号，命题 29）。若 G 单连通，则 \(\rho\) 是基本支配权之和。

\subsection{不可约表示的最高权}

对任意表示 \(\tau : \mathrm{G} \to \mathrm{GL}(\mathrm{V})\)，关联同态 \(\mathrm{L}(\tau)_{(\mathbf{C})}\)，它从 \(\mathbf{C}\)-李代数 \(\mathfrak{g}_{\mathbf{C}}\) 到 \(\operatorname{End}(\mathrm{V})\)，并延拓了线性表示 \(\mathrm{L}(\tau)\)（它是 \(\mathfrak{g}\) 在 V 所对应的实向量空间上的表示）（第 III 章，第 3 节，第 11 号）。根据第 4 节第 3 号的命题 7，映射 \(\delta\) 从 \(\mathrm{X}(\mathrm{T})\) 到 \(\operatorname{Hom}_{\mathbf{C}}(\mathfrak{t}_{\mathbf{C}}, \mathbf{C}) = \mathfrak{t}_{\mathbf{C}}^{*}\) 诱导出一个双射：它将 \(\tau\) 关于 T 的权集合双射到 \(\mathrm{L}(\tau)_{(\mathbf{C})}\) 关于 \(\mathfrak{t}_{\mathbf{C}}\) 的权集合，其中后者是 \(\mathfrak{g}_{\mathbf{C}}\) 的嘉当子代数。

引理 2. 设 \(\varphi\) 是复李代数 \(\mathfrak{g}_{\mathbf{C}}\) 在有限维复向量空间 V 上的线性表示。存在 G 在 V 上的表示 \(\tau\)，使得 \(\mathrm{L}(\tau)_{(\mathbf{C})} = \varphi\)，当且仅当 \(\varphi\) 是半单的且 \(\mathbf{t}_{\mathbf{C}}\) 在 V 上的权属于 \(\delta (\mathrm{X(T)})\)。

若存在 G 的表示 \(\tau\) 使得 \(\mathrm{L}(\tau)_{(\mathbf{C})} = \varphi\)，则 \(\varphi\) 是半单的，因为 G 连通且 \(\tau\) 半单（第 III 章，第 6 节，第 5 号，命题 13 的推论 2），并且 \(\mathfrak{t}_{\mathbf{C}}\) 在 V 上的权属于 \(\delta\) 的像。因此该条件是必要的；下面证明其充分性。若 \(\varphi\) 半单，则 V 是各个 \(V_{\mu}(t_{\mathbf{C}})\) 的直和，其中 \(\mu\) 属于 \(\mathfrak{t}_{\mathbf{C}}\) 在 V 上的权集合（第 VII 章，第 2 节，第 4 号，定理 2 的推论 3）；若所有权都属于 \(\delta\) 的像，则存在 T 在 V 上的表示 \(\tau_{\mathrm{T}}\)，使得 \(\mathrm{L}(\tau_{\mathrm{T}})_{(\mathbf{C})} = \varphi |t_{\mathbf{C}}\)：事实上，只需置 \(\tau_{\mathrm{T}}(t)v = t^{\lambda}v\)，其中 \(t \in \mathrm{T}\) 且 \(v \in V_{\delta(\lambda)}(t_{\mathbf{C}})\)。引理现由第 2 节第 6 号命题 8 得到。

定理 1. a) 设 \(\tau : \mathbf{G} \to \mathbf{GL}(\mathbf{V})\) 是 \(\mathbf{G}\) 的不可约表示。则 \(\tau\) 的权集合（相对于 \(\mathbf{T}\)）有一个最大元 \(\lambda\)，它是支配的，并且空间 \(\mathrm{V}_{\lambda}(\mathrm{T})\) 的维数为 1。

\begin{enumerate}
\def\labelenumi{\alph{enumi})}
\setcounter{enumi}{1}
\item
  G 的两个不可约表示等价，当且仅当它们的最高权相等。
\item
  对于 X(T) 的每个支配元素 \(\lambda\)，存在一个最高权为 \(\lambda\) 的 G 的不可约表示。
\end{enumerate}

由引理 2，G 的不可约表示的等价类与 g 的权属于 \(\delta(\mathrm{X}(\mathrm{T}))\) 的有限维不可约表示的等价类之间存在双射。

将 \(\mathcal{G}\mathfrak{g}_{\mathbf{C}}\) 记为中心，将 \(\mathcal{D}\mathfrak{g}_{\mathbf{C}}\) 记为 \(\mathfrak{g}_{\mathbf{C}}\) 的导出李代数，从而 \(\mathfrak{g}_{\mathbf{C}} = \mathcal{G}\mathfrak{g}_{\mathbf{C}} \oplus \mathcal{D}\mathfrak{g}_{\mathbf{C}}\)。对于每个双线性型 \(\mu\)，其定义域为 \(\mathfrak{t}_{\mathbf{C}} \cap \mathcal{D}\mathfrak{g}_{\mathbf{C}}\)，将第 VIII 章第 6 节第 3 号引入的 \(\mathrm{E}(\mu)\) 记为 \(\mathcal{D}\mathfrak{g}_{\mathbf{C}}\)-单模；对于每个线性形式 \(\nu\)，在 \(\mathcal{G}\mathfrak{g}_{\mathbf{C}}\) 上取与之关联的模 \(\mathbf{C}(\nu)\)，它是 \(\mathcal{G}\mathfrak{g}_{\mathbf{C}}\) 上的一维 \(\mathbf{C}\)-模。则 \(\mathfrak{g}_{\mathbf{C}}\)-模 \(\mathbf{C}(\nu) \otimes \mathrm{E}(\mu)\) 是单模；并且根据第 VIII 章第 7 节第 2 号定理 1 的推论 2 以及《代数》第 VIII 章第 11 节第 1 号定理 1，每个有限维 \(\mathfrak{g}_{\mathbf{C}}\)-单模都同构于这些模之一 \(\mathbf{C}(\nu) \otimes \mathrm{E}(\mu)\)；此外（同上引文），\(\mathbf{C}(\nu) \otimes \mathrm{E}(\mu)\) 是有限维的，当且仅当 \(\mu(H_{\alpha})\) 对每个单根 \(\alpha\) 为正整数。若将 \(\nu + \mu\) 记为 \(\mathfrak{t}_{\mathbf{C}}\) 上的线性形式，它诱导 \(\nu\) 于 \(\mathcal{G}\mathfrak{g}_{\mathbf{C}}\)，并诱导 \(\mu\) 于 \(\mathfrak{t}_{\mathbf{C}} \cap \mathcal{D}\mathfrak{g}_{\mathbf{C}}\)，则 \((\nu + \mu)(H_{\alpha}) = \mu(H_{\alpha})\)；此外，\(\mathbf{C}(\nu) \otimes \mathrm{E}(\mu)\) 的权是 \(\nu + \lambda\)，其中 \(\lambda\) 是 \(\mathrm{E}(\mu)\) 的任一权，因而具有 \(\nu + \mu - \theta\) 的形式，其中 \(\theta \in \delta(X_{+})\)（第 VIII 章，第 6 节，第 2 号，引理 2）。

我们由此得出，\(\mathbf{C}(\nu)\otimes\mathrm{E}(\mu)\) 是有限维的，当且仅当 \((\nu+\mu)(H_{\alpha})\) 对每个单根 \(\alpha\) 为正整数；其权属于 \(\delta(\mathrm{X}(\mathrm{T}))\)，当且仅当 \(\nu+\mu\) 属于 \(\delta(\mathrm{X}(\mathrm{T}))\)。这两个条件同时成立意味着 \(\nu+\mu\) 属于 \(\delta(\mathrm{X}_{++})\)；在此情形下，\(\nu+\mu\) 是 \(\mathbf{C}(\nu)\otimes\mathrm{E}(\mu)\) 的最高权。于是，我们对 X(T) 的每个支配权 \(\lambda\) 都构造出了一个最高权为 \(\lambda\) 的 G 的不可约表示，从而在等价意义下得到了 G 的全部不可约表示。由于权为 \(\nu+\mu\) 的向量构成 \(\mathbf{C}(\nu)\otimes\mathrm{E}(\mu)\) 中的一维子空间，证明完成。

推论. 群 G 有一个（有限维）忠实线性表示。

首先注意，X(T) 的每个元素都等于两个支配权之差；更精确地，令 \(\varpi\) 为 \(X_{++}\) 中满足 \(\langle\varpi,K_{\alpha}\rangle>0\) 对每个单根 \(\alpha\) 成立的元素；对于所有 \(\lambda\in X(T)\)，存在正整数 n，使得 \(\langle\lambda+n\varpi,K_{\alpha}\rangle\geq0\) 对每个单根 \(\alpha\) 成立，即（第 1 号，引理 1）\(\lambda+n\varpi\in X_{++}\)。

因此，存在有限族 \((\lambda_{i})_{i\in I}\) 的 \(X_{++}\) 中元素，生成 Z-模 \(\mathrm{X(T)}\)。对于 \(i\in I\)，令 \(\tau_{i}\) 为最高权为 \(\lambda_{i}\) 的 G 的不可约表示（定理 1）；令表示 \(\tau\) 为各 \(\tau_{i}\) 的直和。按构造，权集 \(\mathrm{P}(\tau,\mathrm{T})\) 是 \(\tau\) 的权集，并生成 Z-模 \(\mathrm{X(T)}\)。由第 4 节第 3 号命题 6，现可知同态 \(\tau\) 是单射，故推论得证。

备注。1) 令 \(\mathfrak{n}_{+}\) 为 \(\mathfrak{g}_{\mathbf{C}}\) 的子代数，它是 \(\mathfrak{g}^{\alpha}\) 对所有 \(\alpha > 0\) 的和。令 \(\tau: \mathrm{G} \to \mathrm{GL}(\mathrm{V})\) 为不可约表示，\(\lambda \in \mathrm{X}_{++}\) 为其最高权，\(\tau': \mathfrak{g}_{\mathbf{C}} \to \mathfrak{gl}(\mathrm{V})\) 为由 \(\tau\) 诱导的表示。则 \(\mathrm{V}_{\lambda}(\mathrm{T})\) 是由向量 \(v\) 属于 \(\mathrm{V}\) 且满足 \(\tau'(x)v = 0\) 对所有 \(x \in \mathfrak{n}_{+}\) 成立的向量组成的子空间。

确实，这由 \(\mathfrak{g}_{\mathbf{C}}\)-模 \(\mathbf{C}(\nu) \otimes \mathrm{E}(\mu)\) 的相应断言得到（第 VIII 章，第 6 节，第 2 号，命题 3）。

\begin{enumerate}
\def\labelenumi{\arabic{enumi})}
\setcounter{enumi}{1}
\item
  令 \(\Theta(\mathrm{G})\) 为 \(\mathrm{G}\) 上取值于 \(\mathbf{C}\) 的连续代表函数的代数（《代数》，第 VIII 章）。令 \(\mathrm{G}\) 通过左、右平移作用于 \(\Theta(\mathrm{G})\)。对于每个 \(\lambda \in \mathrm{X}_{++}\)，令 \((\mathrm{V}_{\lambda}, \tau_{\lambda})\) 为 \(\mathrm{G}\) 的最高权为 \(\lambda\) 的不可约表示（定理 1），令 \((\mathrm{V}_{\lambda}^{*}, \check{\tau}_{\lambda})\) 为反步表示（第 III 章，第 3 节，第 11 号）；根据《谱理论》，\(\mathrm{G} \times \mathrm{G}\) 在 \(\Theta(\mathrm{G})\) 上的表示同构于表示 \((\mathrm{V}_{\lambda} \otimes \mathrm{V}_{\lambda}^{*}, \tau_{\lambda} \otimes \check{\tau}_{\lambda})\) 的直和，其中 \(\lambda\) 遍历 \(\mathrm{X}_{++}\)。备注 1 现在推出如下断言：令 \(\lambda \in \mathrm{X}_{++}\)，令 \(\mathrm{E}_{\lambda}\) 为 \(\Theta(\mathrm{G})\) 的向量子空间，它由连续代表函数 \(f\) 组成，这些函数定义在 \(\mathrm{G}\) 上，满足 \(f(gt) = \lambda(t)^{-1} f(g)\) 对所有 \(g \in \mathrm{G}\) 和所有 \(t \in \mathrm{T}\) 成立，并且满足 \(f * x = 0\) 对所有 \(x \in \mathfrak{n}_{-} = \bigoplus_{\alpha < 0} \mathfrak{g}^{\alpha}\) 成立。则 \(\mathrm{E}_{\lambda}\) 在左平移下稳定，且 \(\mathrm{G}\) 通过左平移在 \(\mathrm{E}_{\lambda}\) 上的表示是最高权为 \(\lambda\) 的不可约表示。
\item
  令 \(\tau : \mathrm{G} \to \mathbf{GL}(\mathrm{V})\) 为不可约表示。存在元素 \(\nu\) 属于 \(\mathrm{X}(\mathrm{C}(\mathrm{G}))\)，使得 \(\tau(s)v = \nu(s)v\) 对所有 \(s \in \mathrm{C}(\mathrm{G}), v \in \mathrm{V}\) 成立：事实上，\(\tau(\mathrm{C}(\mathrm{G}))\) 包含于 \(\tau(\mathrm{G})\) 的交换子中，而后者等于 \(\mathbf{C}^*.1_{\mathrm{V}}\)（《代数》，第 VIII 章，第 3 节，第 2 号，定理 1）。对于每个权 \(\lambda\) 属于 \(\tau\)，\(\lambda\) 在 \(\mathrm{C}(\mathrm{G})\) 上的限制等于 \(\nu\)。
\item
  第 VIII 章第 7 节第 2 号至第 5 号的定义和断言可以毫无困难地推广到当前情形；细节留给读者。
\end{enumerate}

命题 1. 设 \(\tau : G \to \mathbf{GL}(V)\) 是 \(G\) 的最高权为 \(\lambda \in X_{++}\) 的不可约表示。令 \(m\) 为整数 \(\sum_{\alpha \in R_+} \langle \lambda, K_\alpha \rangle\)，令 \(w_0\) 为满足 \(w_0(R_+) = R_-\) 的 Weyl 群元素（第 VI 章，第 6 节，命题 17 的推论 3）。有三种可能的情形：

\begin{enumerate}
\def\labelenumi{\alph{enumi})}
\item
  \(w_{0}(\lambda) = -\lambda\) 且 m 为偶数。则 V 上存在 G-不变的分离对称双线性型；表示 \(\tau\) 为实型（附录 II）。
\item
  \(w_{0}(\lambda) \neq -\lambda\)。V 上每个 G-不变双线性型都为零；表示 \(\tau\) 为复型（同上引文）。
\item
  \(w_{0}(\lambda) = -\lambda\) 且 m 为奇数。则 V 上存在 G-不变的分离交错双线性型；表示 \(\tau\) 为四元数型（同上引文）。
\end{enumerate}

若 \(\tau\) 在 \(\mathrm{C}(\mathrm{G})_0\) 上的限制不是平凡的，则属于情形 \(b\)。

V 上的双线性型 B 在 G 下不变，当且仅当它在 \(g_{C}\) 下不变（第 III 章，第 6 节，第 5 号，推论 3）。因此，若 G 为半单的，则命题由第 VIII 章第 7 节第 5 号命题 12 以及附录 II 的命题 3 得出。

在一般情形下，置 \(\mathrm{C}(\mathrm{G})_{0} = \mathrm{S}\)，并将 \(\mathrm{X}(\mathrm{T}/\mathrm{S})\) 与 \(\mathrm{X}(\mathrm{T})\) 的一个子群识别（该子群在 W 下稳定）。若 \(\tau(\mathrm{S}) = \{1_{\mathrm{V}}\}\)，则 \(\tau\) 通过取商诱导出 G/S 的表示 \(\tau' : G/S \to GL(V)\)，其最高权为 \(\lambda\)；在此情形下，将前面的结果应用于 G/S 即得命题。

假设 \(\tau (\mathrm{S})\neq \{1_{\mathrm{V}}\}\)。存在非零元素 \(\nu\) 属于 \(\mathrm{X(S)}\)，使得 \(\tau (s) = \nu (s)_{\mathrm{V}}\) 对所有 \(s\in \mathrm{S}\) 成立（备注 3）。于是 \(\nu\) 是 \(\lambda\) 在限制同态 \(\mathrm{X}(\mathrm{T})\to\mathrm{X}(\mathrm{S})\) 下的像；由于 W 在 \(\mathrm{X}(\mathrm{S})\) 上平凡作用，等式 \(w_{0}(\lambda)=-\lambda\) 蕴含 \(\nu=-\nu\)，这是不可能的；所以 \(w_{0}(\lambda)\neq-\lambda\)。另一方面，若 B 是 V 上 G-不变的双线性型，则对于 V 中所有 x,y 以及 S 中 s，有

\[
\mathrm{B} (\nu (s) x, \nu (s) y) = \mathrm{B} (x, y) = \nu (s) ^ {2} \mathrm{B} (x, y)
\]

这蕴含 B = 0，故命题得证。

令 \(\mathfrak{S}_{\mathbf{R}}(\mathrm{G})\) 为 G 在有限维实向量空间上的不可约连续表示的等价类集合。命题 1 及附录 II 的结果给出双射 \(\Phi:X_{++}/\Sigma\to\mathfrak{S}_{\mathbf{R}}(\mathrm{G})\)，其中 \(\Sigma\) 表示子群 \(\{1,-w_{0}\}\) \(\operatorname{Aut}(\mathrm{X}(\mathrm{T}))\) 的子群。更精确地，令 \(\lambda\in X_{++}\)，令 \(E_{\lambda}\) 为最高权为 \(\lambda\) 的 G 的一个表示；则

\[
\begin{array}{c} \Phi (\{\lambda , - w _ {0} (\lambda) \}) = \mathrm{E} _ {\lambda [ \mathbf {R} ]} \text {if} \lambda \neq - w _ {0} (\lambda) \text {or if} \sum_ {\alpha \in \mathrm{R} _ {+}} \langle \lambda , K _ {\alpha} \rangle \notin 2 \mathbf {Z} \\ \Phi (\{\lambda , - w _ {0} (\lambda) \}) = \mathrm{E} _ {\lambda} ^ {\prime} \text {if} \lambda = - w _ {0} (\lambda) \text {and} \sum_ {\alpha \in \mathrm{R} _ {+}} \langle \lambda , K _ {\alpha} \rangle \in 2 \mathbf {Z} \end{array}
\]

其中 \(E_{\lambda}^{\prime}\) 是 \(E_{\lambda}\) 上在 G 下不变的 R-结构。

\subsection{表示环 R(G)}

令 \(\mathrm{R}(\mathrm{G})\) 为 \(\mathrm{G}\) 的表示环（即有限维复向量空间上的连续表示）（《代数》，第 VIII 章，第 10 节，第 6 号）。若 \(\tau\) 是 \(\mathrm{G}\) 的表示，将 \([\tau]\) 记为 \(\mathrm{R}(\mathrm{G})\) 中的等价类；若 \(\tau\) 和 \(\tau'\) 是 \(\mathrm{G}\) 的两个表示，则按定义有

\[
\begin{array}{r} [ \tau ] + [ \tau^ {\prime} ] = [ \tau \oplus \tau^ {\prime} ] \\ \relax [ \tau ] [ \tau^ {\prime} ] = [ \tau \otimes \tau^ {\prime} ]. \end{array}
\]

由于 G 的每个表示都是半单的，Z-模 \(\mathrm{R}(\mathrm{G})\) 是自由的，并以 G 的不可约表示的等价类集合为基；根据定理 1，该集合可与 \(X_{++}\) 识别。映射 \(\tau \mapsto \mathrm{L}(\tau)_{(\mathbf{C})}\) 诱导出从环 \(\mathrm{R}(\mathrm{G})\) 到环 \(\mathcal{R}(\mathfrak{g}_{\mathbf{C}})\) 的同态 l，它是 \(g_{C}\) 的表示环（第 VIII 章，第 7 节，第 6 号）。

令 \(\tau : \mathbf{G} \to \mathbf{GL}(\mathbf{V})\) 为 \(\mathbf{G}\) 的表示；考虑分次 \((\mathrm{V}_{\lambda}(\mathrm{T}))_{\lambda \in \mathrm{X}(\mathrm{T})}\)，它是 \(\mathbf{C}\)-向量空间 \(\mathbf{V}\) 的分次。将 \(\mathrm{Ch}(\mathbf{V})\) 或 \(\mathrm{Ch}(\tau)\) 记为分次向量空间 \(\mathbf{V}\) 的特征标（第 VIII 章，第 7 节，第 7 号）；若 \((e^{\lambda})_{\lambda \in \mathrm{X}(\mathrm{T})}\) 表示代数 \(\mathbf{Z}[\mathrm{X}(\mathrm{T})] = \mathbf{Z}^{(\mathrm{X}(\mathrm{T}))}\) 的典则基，则按定义有

\[
\operatorname{Ch} (\tau) = \sum_ {\lambda \in \mathrm{X} (\mathrm{T})} \left(\dim \mathrm{V} _ {\lambda} (\mathrm{T})\right) e ^ {\lambda}.
\]

以同样方式（同上引文）定义从 R(G) 到 Z[X(T)] 的环同态，仍记为 Ch。若 G 是半单的，则有交换图

\[
\begin{array}{c c c} \mathrm{R(G)} & \stackrel {{\mathrm{Ch}}} {{\longrightarrow}} & \mathbf {Z} [ \mathrm{X(T)} ] \\ \Big \downarrow_ {l} & & \Big \downarrow_ {\tilde {\delta}} \\ \mathcal {R} (\mathfrak {g} _ {\mathbf {C}}) & \stackrel {{\mathrm{ch}}} {{\longrightarrow}} & \mathbf {Z} [ \mathrm{P} ] \end{array}\tag{1}
\]

其中 P 表示 \(\mathrm{R}(\mathfrak{g}_{\mathbf{C}}, \mathfrak{t}_{\mathbf{C}})\) 的权群，\(\tilde{\delta}\) 表示由 \(\delta\) 诱导的同态。

Weyl 群 W 通过自同构作用于群 X(T)，因而作用于环 Z[X(T)]。根据第 4 节第 3 号命题 5，Ch 的像包含于子环 Z[X(T)] \(^{W}\) 中，该子环由在 W 下不变的元素组成。

命题 2. 同态 \(\mathrm{Ch}\) 诱导出从 \(\mathrm{R}(\mathrm{G})\) 到 \(\mathrm{Z}[\mathrm{X}(\mathrm{T})]^{\mathrm{W}}\) 的同构。

对于 \(\lambda \in \mathrm{X}_{++}\)，将 \([\lambda]\) 记为 \(\mathrm{R}(\mathrm{G})\) 中最高权为 \(\lambda\) 的不可约表示的等价类。由于族 \(([\lambda])_{\lambda \in \mathrm{X}_{++}}\) 是 \(\mathbf{Z}\)-模 \(\mathrm{R}(\mathrm{G})\) 的一个基，只需证明以下断言：

族 \((\mathrm{Ch}[\lambda])_{\lambda\in\mathrm{X}_{++}}\) 是 Z-模 \(\mathbf{Z}[\mathrm{X}(\mathrm{T})]^{\mathrm{W}}\) 的一个基。

对于元素 \(u = \sum_{\lambda} a_{\lambda} e^{\lambda}\)（属于 \(\mathbf{Z}[X(T)]\)），其中的项 \(a_{\lambda} e^{\lambda}\)，若 \(\lambda\) 是由 \(\mu \in X(T)\) 且 \(a_{\mu} \neq 0\) 组成的集合的最大元，则称该项为极大项。定理 1 蕴含 \(\operatorname{Ch}[\lambda]\) 有唯一的极大项，即 \(e^{\lambda}\)。因此，命题由下述引理得到。

引理 3. 对每个 \(\lambda \in \mathrm{X}_{++}\)，令 \(\mathrm{C}_{\lambda}\) 为 \(\mathbf{Z}[\mathrm{X(T)}]^{\mathrm{W}}\) 中具有唯一极大项 \(e^{\lambda}\) 的元素。则族 \((\mathrm{C}_{\lambda})_{\lambda \in \mathrm{X}_{++}}\) 是 \(\mathbf{Z}[\mathrm{X(T)}]^{\mathrm{W}}\) 的一个基。

证明与第 VI 章第 3 节第 4 号命题 3 的证明相同，只需将 A 换成 Z，将 P 换成 X(T)，并将 P ∩ \(\overline{C}\) 换成 \(X_{++}\)。

令 \(\Theta (\mathrm{G})\)（分别为 \(\Theta (\mathrm{T})\)）为取值于 \(\mathbf{C}\) 的 \(\mathrm{G}\)（分别为 T）上连续代表函数的代数，并令 \(\mathrm{Z}\Theta (\mathrm{G})\)（分别为 \(\Theta (\mathrm{T})^{\mathrm{W}}\)）为由中心函数（分别为 W-不变函数）组成的子代数。限制映射 \(\Theta (\mathrm{G})\to \Theta (\mathrm{T})\) 诱导出环同态 \(r:\mathrm{Z}\Theta (\mathrm{G})\rightarrow \Theta (\mathrm{T})^{\mathrm{W}}\)。另一方面，将表示 \(\tau\) 与其特征标对应的映射（即函数 \(g\mapsto \mathrm{Tr}\tau (g)\)）延拓为 \(\mathbf{C}\)-代数同态 \(\mathrm{Tr}:\mathbf{C}\otimes_{\mathbf{Z}}\mathrm{R}(\mathrm{G})\to \mathrm{Z}\Theta (\mathrm{G})\)；根据《谱理论》，它是同构。同样，典则嵌入 \(\mathrm{X(T)}\to \Theta (\mathrm{T})\) 诱导出 \(\mathbf{C}\)-代数同构 \(\iota :\mathbf{C}[\mathrm{X(T)}]\to \Theta (\mathrm{T})\)，进而诱导同构 \(\iota :\mathbf{C}[\mathrm{X(T)}]^{\mathrm{W}}\to \Theta (\mathrm{T})^{\mathrm{W}}\)。图

\[
\begin{array}{c c c} \mathbf {C} \otimes_ {\mathbf {Z}} \mathrm{R(G)} & \stackrel {{1 \otimes \mathrm{Ch}}} {{\longrightarrow}} & \mathbf {C} [ \mathrm{X(T)} ] ^ {\mathrm{W}} \\ \Big \downarrow_ {\mathrm{Tr}} & & \Big \downarrow_ {\iota} \\ \mathrm{Z} \Theta (\mathrm{G}) & \stackrel {{r}} {{\longrightarrow}} & \Theta (\mathrm{T}) ^ {\mathrm{W}} \end{array}\tag{2}
\]

是交换的：事实上，对于每个表示 \(\tau : \mathrm{G} \to \mathbf{GL}(\mathrm{V})\) 以及所有 \(t \in \mathrm{T}\)，

\[
\operatorname{Tr} \tau (t) = \sum_ {\lambda \in \mathrm{X} (\mathrm{T})} \left(\dim \mathrm{V} _ {\lambda} (\mathrm{T})\right) \lambda (t) = \iota (\operatorname{Ch} \tau) (t),
\]

即 \((r\circ \mathrm{Tr})[\tau ] = (\iota \circ \mathrm{Ch})[\tau ]\)

现可由命题推出以下结果。

推论. 限制映射 \(r: \mathrm{Z}\Theta(\mathrm{G}) \to \Theta(\mathrm{T})^{\mathrm{W}}\) 是双射。

\subsection{特征标公式}

在本号中，将群 \(X(T)\) 写成乘法记号，并将其元素视为 T 上的复值函数。假设元素 \(\rho\) 是 \(X(T) \otimes Q\) 中属于 \(X(T)\) 的元素。

将 \(L^{2}(T)\) 记为 T 上平方可积复函数的等价类组成的 Hilbert 空间，将 \(\Theta(T)\) 记为连续代表函数组成的子空间。根据《谱理论》，\(X(T)\) 是 \(L^{2}(T)\) 的正交归一基，也是 \(\Theta(T)\) 的代数基。

对于 \(f \in \mathrm{L}^{2}(\mathrm{T})\) 和 \(w \in W\)，将 \(\mathrm{L}^{2}(\mathrm{T})\) 中由 \(wf(t) = f(w^{-1}(t))\) 定义的元素 wf 记为；于是，对于 \(\lambda \in \mathrm{X}(\mathrm{T})\)，有 \(w\lambda = w(\lambda)\)。将 \(\varepsilon : W \to \{1, -1\}\) 记为符号（即满足对每个反射 s 都有 \(\varepsilon(s) = -1\) 的唯一同态）；对于 \(f \in \mathrm{L}^{2}(\mathrm{T})\)，置

\[
\mathrm{J} (f) = \sum_ {w \in \mathrm{W}} \varepsilon (w) ^ {w} f.
\]

若 \(\lambda \in X_{++}\)，则各特征标 \(^w (\lambda \rho)\) 彼此不同；事实上，只需证明 \(^w (\lambda \rho)\neq \lambda \rho\) 对所有 \(w\neq 1\) 成立；但这由引理 1（第 1 号）以及事实 \(\langle \lambda \rho ,K_{\alpha}\rangle = \langle \lambda ,K_{\alpha}\rangle +1 > 0\) 对每个正根 \(\alpha\) 成立得到。因此，

\[
\| \mathrm{J} (\lambda \rho) \| ^ {2} = \operatorname{Card} (\mathrm{W}) = w (\mathrm{G}).
\]

称 \(f \in \mathrm{L}^{2}(\mathrm{T})\) 为反不变元，若 \(^{w}f = \varepsilon(w)f\) 对所有 \(w \in W\) 成立（即，对每个反射 s 都有 \(^{s}f = -f\)）。我们将证明 \(\frac{1}{w(\mathrm{G})}\mathrm{J}\) 是 \(\mathrm{L}^{2}(\mathrm{T})\) 到反不变元子空间上的正交投影。事实上，令 \(f, f'\) 属于 \(\mathrm{L}^{2}(\mathrm{T})\)，且 \(f'\) 反不变；则 \(\mathrm{J}(f)\) 反不变，并且

\[
\begin{array}{c} \langle f ^ {\prime}, \mathrm{J} (f) \rangle = \sum_ {w \in \mathrm{W}} \varepsilon (w) \langle f ^ {\prime},   ^ {w} f \rangle = \sum_ {w \in \mathrm{W}} \langle^ {w} f ^ {\prime},   ^ {w} f \rangle \\ = \sum_ {w \in \mathrm{W}} \langle f ^ {\prime},   f \rangle = w (\mathrm{G}) \langle f ^ {\prime},   f \rangle . \end{array}
\]

命题 3. 元素 \(\mathrm{J}(\lambda\rho)/\sqrt{w(\mathrm{G})}\)，其中 \(\lambda\in X_{++}\)，构成 \(\mathrm{L}^{2}(\mathrm{T})\) 的反不变元子空间的正交归一基，并构成 \(\Theta(\mathrm{T})\) 的反不变元子空间的代数基。

证明与第 VI 章第 3 节第 3 号命题 1 的证明相同。

根据第 VI 章同上引文的命题 2，

\[
\mathrm{J} (\rho) = \rho \prod_ {\alpha > 0} (1 - \alpha^ {- 1}) = \rho^ {- 1} \prod_ {\alpha > 0} (\alpha - 1),\tag{3}
\]

因此

\[
\mathrm{J} (\rho) \overline {{\mathrm{J} (\rho)}} = \prod_ {\alpha} (\alpha - 1).\tag{4}
\]

由定理 1 的推论 2（第 6 节，第 2 号），得到：

引理 4. 若 \(\varphi\) 和 \(\psi\) 是 \(\mathbf{G}\) 上的两个连续中心函数，

\[
\int_ {\mathrm{G}} \overline {{\varphi (g)}} \psi (g) d g = \frac {1}{w (\mathrm{G})} \int_ {\mathrm{T}} \overline {{(\varphi (t) \mathrm{J} (\rho) (t))}}. (\psi (t) \mathrm{J} (\rho) (t)) d t.
\]

对于所有 \(\lambda\in X_{++}\)，将 \(\chi_{\lambda}\) 记为最高权为 \(\lambda\) 的 G 的不可约表示的特征标。

定理 2（H. Weyl）。对于所有 \(\lambda \in \mathrm{X}_{++}\)，有 \(\mathrm{J}(\rho)\chi_{\lambda}|\mathrm{T} = \mathrm{J}(\lambda \rho)\)。

函数 \(\mathrm{J}(\rho)\cdot\chi_{\lambda}|T\) 在 W 下反不变，并且是 \(\mathrm{X(T)}\) 中元素的整数系数线性组合。因此，根据第 VI 章第 3 节第 3 号命题 1，它可以写成 \(\sum_{\mu}a_{\mu}\mathrm{J}(\mu\rho)\)，其中 \(\mu\) 属于 \(X_{++}\)，且 \(a_{\mu}\) 是整数，除有限多个外均为零；由于 \(\int_{\mathrm{G}}|\chi_{\lambda}(g)|^{2}dg=1\)

（《谱理论》），由命题 3 和引理 4 得到 \(\sum_{\mu}(a_{\mu})^{2} = 1\)；因此 \(a_{\mu}\) 除一个外全为零，而该值必为 1 或 \(-1\)。但是，\(\lambda\) 在 \(\chi_{\alpha}|T\) 中的系数等于 1（定理 1），所以 \(\lambda \rho\) 在 \(J(\rho)\cdot \chi_{\lambda}|T\) 中的系数等于 1（第 VI 章，第 3 节，第 3 号，备注 2），这意味着 \(a_{\lambda \rho} = 1\)，定理得证。

推论 1. 按第 3 号的记号，在 \(\mathbf{Z}[\mathrm{X(T)}]\) 中有

\[
\left(\sum_ {w \in \mathrm{W}} \varepsilon (w) e ^ {w \rho}\right) \operatorname{Ch} [ \lambda ] = \sum_ {w \in \mathrm{W}} \varepsilon (w) e ^ {w \lambda} e ^ {w \rho} \text {for all} \lambda \in \mathrm{X} _ {+ +}.
\]

这由定理和交换图（2）（第 3 号）得到。

推论 2. 对所有 \(\lambda \in X_{++}\) 以及 T 的每个正则元 \(t\)，

\[
\chi_ {\lambda} (t) = \frac {\sum_ {w} \varepsilon (w) \lambda (w t) \rho (w t)}{\sum_ {w} \varepsilon (w) \rho (w t)}\tag{5}
\]

其中两个求和遍历 W 的元素 w。

事实上，\(\mathrm{J}(\rho)(t)\) 非零，因为 \(t\) 是正则元（公式（4））。

若 \(\varphi\) 是 \(\mathbf{G}\) 上的中心函数，则 \(\varphi\) 在 \(\mathrm{T}\) 上的限制在 \(\mathrm{W}\) 下不变，所以 \(\mathrm{J}(\rho). \varphi|\mathrm{T}\) 在 \(\mathrm{W}\) 下反不变。此外，根据《谱理论》和定理 1，族 \((\chi_{\lambda})_{\lambda \in \mathrm{X}_{++}}\) 是 \(\mathbf{G}\) 上中心代表函数空间的代数基，也是 \(\mathrm{ZL}^2(\mathrm{G})\)（即 \(\mathbf{G}\) 上平方可积中心函数的等价类空间）的正交归一基。

因此，由命题 3 和定理 2 得到：

推论 3. 将 G 上任意连续中心函数 \(\varphi\) 映为函数 \(w(\mathrm{G})^{-1/2}\mathrm{J}(\rho)(\varphi|\mathrm{T})\) 的映射，诱导出从 G 上中心代表函数空间到 \(\Theta(\mathrm{T})\) 的反不变元空间的同构；它由连续性延拓为从 \(\mathrm{ZL}^2(\mathrm{G})\) 到 \(\mathrm{L}^2(\mathrm{T})\) 的反不变元子空间的 Hilbert 空间同构。

推论 4. 令 \(\varphi\) 为 G 上的连续中心函数。则

\[
\int_ {\mathrm{G}} \overline {{\chi_ {\lambda} (g)}} \varphi (g) d g = \int_ {\mathrm{T}} \overline {{\lambda (t)}} \prod_ {\alpha > 0} (1 - \alpha (t) ^ {- 1}) \varphi (t) d t = \int_ {\mathrm{T}} \overline {{\lambda \rho (t)}}. \varphi (t) \mathrm{J} (\rho) (t) d t.
\]

事实上，由引理 4 和定理 2，

\[
\begin{array}{c} \int_ {\mathrm{G}} \overline {{\chi_ {\lambda} (g)}} \varphi (g) d g = \frac {1}{w (\mathrm{G})} \int_ {\mathrm{T}} \overline {{\chi_ {\lambda} (t) \mathrm{J} (\rho) (t)}} (\varphi (t) \mathrm{J} (\rho) (t)) d t \\ = \frac {1}{w (\mathrm{G})} \int_ {\mathrm{T}} \overline {{\mathrm{J} (\lambda \rho) (t)}} \varphi (t) \mathrm{J} (\rho) (t) d t. \end{array}
\]

但是函数 \(t \mapsto \varphi(t)J(\rho)(t)\) 是反不变的，而 \(\frac{1}{w(G)}J(\lambda\rho)\) 是将 \(\lambda\rho\) 投影到 \(L^{2}(T)\) 的反不变元子空间上的正交投影，因此

\[
\frac {1}{w (\mathrm{G})} \int_ {\mathrm{T}} \overline {{\mathrm{J} (\lambda \rho) (t)}} \varphi (t) \mathrm{J} (\rho) (t) d t = \int_ {\mathrm{T}} \overline {{\lambda \rho (t)}} \varphi (t) \mathrm{J} (\rho) (t) d t;
\]

最后，由公式（3），有 \(\overline{\rho(t)}\mathrm{J}(\rho)(t) = \prod_{\alpha >0}(1 - \alpha (t)^{-1})\)，故推论得证。

备注。1) 对于所有 \(w \in \mathrm{W}\)，置 \(\rho_w = {}^w\rho /\rho\)；则

\[
\sum_ {w} \varepsilon (w) \rho_ {w} = \prod_ {\alpha > 0} (1 - \alpha^ {- 1}) = \rho^ {- 2} \prod_ {\alpha > 0} (\alpha - 1).\tag{6}
\]

若 \(t\) 是 \(\mathrm{T}\) 的正则元，则由（5）得到

\[
\chi_ {\lambda} (t) = \frac {\sum_ {w} \varepsilon (w) ^ {w} \lambda (t) \rho_ {w} (t)}{\sum_ {w} \varepsilon (w) \rho_ {w} (t)} = \frac {\sum_ {w} \varepsilon (w) ^ {w} \lambda (t) \rho_ {w} (t)}{\prod_ {\alpha > 0} (1 - \alpha (t) ^ {- 1})}.\tag{7}
\]

注意，\(\rho_{w}\) 是根的整数系数线性组合，因而即使不假设 \(\rho \in \mathrm{X}(T)\)，它仍属于 X(T)。由此可知，公式（7）无需假设 \(\rho \in \mathrm{X}(T)\) 也成立：事实上，为证明这一点，只需将 G 换成适当的连通覆盖，便可归结为推论 2。

\begin{enumerate}
\def\labelenumi{\arabic{enumi})}
\setcounter{enumi}{1}
\item
  同样，推论 4 的第一个等式无需假设 \(\rho \in \mathrm{X}(\mathrm{T})\) 也成立。
\item
  可以由相应的无穷小断言（第 VIII 章，第 9 节，第 1 号，定理 1）推出定理 2；下一号的定理 3 也同样如此（它是上述引文第 2 号定理 2 的类似结果）。
\end{enumerate}

\subsection{不可约表示的次数}

现在回到群 \(X(T)\) 的加法记号，并且不再假设 \(\rho\) 属于 \(X(T)\)。

定理 3. 最高权为 \(\lambda\) 的 G 的不可约表示空间的维数由下式给出

\[
\chi_ {\lambda} (e) = \prod_ {\alpha \in \mathrm{R} _ {+}} \frac {\langle \lambda + \rho , K _ {\alpha} \rangle}{\langle \rho , K _ {\alpha} \rangle}.
\]

置 \(\gamma = \frac{1}{2}\sum_{\alpha >0}K_{\alpha}\)，于是 \(\delta (\alpha)(\gamma) = 2\pi i\) 对每个单根 \(\alpha\) 成立（第 VI 章，第 1 节，第 10 号，命题 29）。直线 \(\mathbf{R}\gamma\) 不包含在任何超平面 \(\mathrm{Ker}\delta (\alpha)\) 中，所以 \(\exp (z\gamma)\) 对所有充分小的 \(z\in \mathbf{R}^*\) 都是正则元；对于所有 \(\mu \in \mathrm{X(T)}\) 和所有 \(z\in \mathbf{R}\)，有

\[
\mathrm{J} (\mu) (\exp (z \gamma)) = \sum_ {w \in \mathrm{W}} \varepsilon (w) e ^ {z \delta (\mu) (w ^ {- 1} \gamma)}.
\]

暂且接受下述引理：

引理 5. 有

\[
\mathrm{J} (\mu) (\exp (z \gamma)) = e ^ {z \delta (\mu) (\gamma)} \prod_ {\alpha > 0} (1 - e ^ {- z \delta (\mu) (K _ {\alpha})}).
\]

因此，函数 \(\mathrm{J}(\mu)(\exp (z\gamma))\) 是一个当 \(z\) 趋于 0 时趋于 1 的函数与下式的乘积

\[
z ^ {\mathrm{N}} \prod_ {\alpha > 0} \delta (\mu) (K _ {\alpha}) = (2 \pi i z) ^ {\mathrm{N}} \prod_ {\alpha > 0} \langle \mu , K _ {\alpha} \rangle
\]

其中 N = R+ 的基数。

首先假设 \(\rho \in \mathrm{X(T)}\)；应用定理 2 的推论 2，可知当 \(z\) 趋于 0 时，\(\chi_{\lambda}(z\gamma)\) 趋于

\[
\prod_ {\alpha > 0} \langle \lambda + \rho , K _ {\alpha} \rangle / \prod_ {\alpha > 0} \langle \rho , K _ {\alpha} \rangle ,
\]

因此在此情形下定理成立。

在一般情形下，只需注意，在证明定理 3 时总可以将 G 换成适当的连通覆盖，从而归结为前面的情形。

现在证明引理 5。令 \(z \in \mathbf{C}\)；将 \(\varphi_z\) 记为从 \(\mathfrak{t}\) 到 \(\mathbf{C}\)-代数 \(\operatorname{Map}(\mathrm{X}(\mathrm{T}), \mathbf{C})\) 的映射；这里的代数由从 \(\mathrm{X}(\mathrm{T})\) 到 \(\mathbf{C}\) 的映射组成；它将 \(H \in \mathfrak{t}\) 对应到映射

\[
\varphi_ {z} (H): \mu \mapsto \mu (\exp z H) = e ^ {z \delta (\mu) (H)}.
\]

有 \(\varphi_z(H + H') = \varphi_z(H)\varphi_z(H')\)，所以存在环同态

\[
\psi_ {z}: \mathbf {Z} [ \mathfrak {t} ] \to \operatorname{Map} (\mathrm{X} (\mathrm{T}), \mathbf {C})
\]

使得 \(\psi_z(e^H)(\mu) = e^{z\delta (\mu)(H)}\)。另一方面，根据第 VI 章第 3 节第 3 号命题 2，在 \(\mathbf{Z}[t]\) 中有如下关系：

\[
\sum_ {w \in \mathrm{W}} \varepsilon (w) e ^ {w \gamma} = e ^ {\gamma} \prod_ {\alpha > 0} (1 - e ^ {- K _ {\alpha}}).
\]

应用同态 \(\psi_z\)，并利用等式

\[
\psi_ {z} (e ^ {w \gamma}) (\mu) = e ^ {z \delta (\mu) (w \gamma)} = e ^ {z \delta (w ^ {- 1} \mu) (\gamma)},
\]

得到所述公式。

推论 1. 令 \(\| \|\) 为 \(\mathrm{X(T)}\otimes \mathbf{R}\) 上的范数。对于所有 \(\lambda \in \mathrm{X}_{++}\)，令 \(d(\lambda)\) 为 \(\mathrm{G}\) 的最高权为 \(\lambda\) 的不可约表示空间的维数。

\(a) \sup_{\lambda \in X_{++}} d(\lambda) / \| \lambda + \rho \|^{N} < \infty, \text{ 其中 } N = 1/2 (\dim G - \dim T).\)

\begin{enumerate}
\def\labelenumi{\alph{enumi})}
\setcounter{enumi}{1}
\item
  若 G 是半单的，则 \(\inf_{\lambda\in X_{++}}d(\lambda)/\|\lambda+\rho\|>0\)。
\item
  对于所有 \(\alpha \in \mathrm{R}_{+}\)，存在 \(\mathrm{A}_{\alpha} > 0\)，使得 \(|\langle \lambda + \rho, K_{\alpha} \rangle| \leq \mathrm{A}_{\alpha} \| \lambda + \rho \|\)，因而 \(d(\lambda) / \| \lambda + \rho \|^{\mathrm{N}} \leq \prod_{\alpha > 0} \mathrm{A}_{\alpha} / \langle \rho, K_{\alpha} \rangle\)。
\item
  假设 G 是半单的，将单根记为 \(\beta_{1},\ldots,\beta_{r}\)，并置 \(N_{i}=K_{\beta_{i}}\)。则
\end{enumerate}

\[
d (\lambda) \geq \prod_ {i = 1} ^ {r} \frac {\langle \lambda + \rho , N _ {i} \rangle}{\langle \rho , N _ {i} \rangle} = \prod_ {i = 1} ^ {r} \langle \lambda + \rho , N _ {i} \rangle ;
\]

由于 \(\langle \lambda +\rho ,N_i\rangle \geq \langle \rho ,N_i\rangle = 1\)，这意味着

\[
d (\lambda) \geq \sup _ {i} | \langle \lambda + \rho , N _ {i} \rangle |.
\]

若 G 是半单的，则 \(x \mapsto \sup |\langle x, N_{i} \rangle|\) 是 \(\mathrm{X(T)} \otimes \mathbf{R}\) 上的范数，必然与给定范数等价，故 b) 成立。

推论 2. 假设 G 是半单的；令 d 为整数。维数 \(\leq d\) 的 G 的表示的等价类集合是有限的。

推论 1 b) 蕴含，集合 \(X_{d}\) 中的元素 \(\lambda\) 属于 \(X_{++}\) 且满足 \(d(\lambda) \leq d\) 的元素所成的集合是有限的。对于所有 \(\lambda\) 属于 \(X_{d}\)，令 \(V_{\lambda}\) 为最高权为 \(\lambda\) 的不可约表示；每个维数 \(\leq d\) 的表示都同构于直和 \(\bigoplus_{\lambda \in X_{d}} V_{\lambda}^{n_{\lambda}}\)，其中 \(n_{\lambda} \leq d\)，故推论得证。

\subsection{CASIMIR 元}

根据第 1 节第 3 号命题 3，g 上存在负的、对称的、分离的且在 Ad(G) 下不变的双线性型（例如，若 G 是半单的，可以取 g 的 Killing 型）。令 F 为这样的一个型。回忆（第 I 章，第 3 节，第 7 号），与 F 关联的 Casimir 元是包络代数 U(g) 的中心元素 \(\Gamma\)，满足：对于 g 的任意基 \((e_{i})\)，若 \(\mathrm{F}(e_{i}, e_{j}) = -\delta_{ij}\)，则有 \(\Gamma = -\sum e_{i}^{2}\)。

在本章余下部分中，我们把 G 的 Casimir 元称为由 g 上负的、分离的、不变的对称双线性型按上述方式得到的 \(\mathrm{U}(\mathfrak{g})\) 的元素。若 \(\Gamma\) 是 G 的 Casimir 元，且 \(\tau: G \to \mathbf{GL}(V)\) 是 G 的不可约表示，则 V 的自同态 \(\Gamma_{V}\) 是一个伸缩变换（《代数》，第 VIII 章，第 3 节，第 2 号，定理 1），其比率记为 \(\tilde{\Gamma}(\tau)\)。

命题 4. 令 \(\Gamma\) 为 G 的 Casimir 元。

\begin{enumerate}
\def\labelenumi{\alph{enumi})}
\item
  若 \(\tau\) 是 \(\mathbf{G}\) 的不可约表示，则 \(\tilde{\Gamma}(\tau)\) 为正实数。若 \(\tau\) 不是平凡表示，则 \(\tilde{\Gamma}(\tau) > 0\)。
\item
  存在唯一的二次型 \(\mathbf{Q}_{\Gamma}\)，定义在 \(\mathrm{X(T)}\otimes \mathbf{R}\) 上，使得对于每个不可约表示 \(\tau\)，它属于 \(\mathbf{G}\)，
\end{enumerate}

\[
\tilde {\Gamma} (\tau) = \mathrm{Q} _ {\Gamma} (\lambda + \rho) - \mathrm{Q} _ {\Gamma} (\rho),
\]

其中 \(\lambda\) 是 \(\tau\) 的最高权。型 \(Q_{\Gamma}\) 是正的、分离的且在 W 下不变。

令 \(\mathrm{F}\) 为 \(\mathfrak{g}\) 上定义 \(\Gamma\) 的分离负对称双线性型。令 \(\tau : \mathrm{G} \to \mathbf{GL}(\mathrm{V})\) 为 \(\mathrm{G}\) 的不可约表示；令 \(\langle , \rangle\) 为 \(\mathrm{V}\) 上在 \(\mathrm{G}\) 下不变的 Hilbert 内积（第 1 节，第 1 号），并令 \((e_i)\) 为 \(\mathfrak{g}\) 的满足 \(\mathrm{F}(e_i, e_j) = -\delta_{ij}\) 的基。则对于每个元素 \(v\)，它属于 \(\mathrm{V}\) 且不在 \(\mathrm{G}\) 下不变，有

\[
\begin{array}{c} \tilde {\Gamma} (\tau) \langle v, v \rangle = \langle v, \Gamma_ {\mathrm{V}} (v) \rangle = - \sum_ {i} \langle v, (e _ {i}) _ {\mathrm{V}} ^ {2} v \rangle = \sum_ {i} \langle v, (e _ {i}) _ {\mathrm{V}} ^ {*} (e _ {i}) _ {\mathrm{V}} v \rangle \\ = \sum_ {i} \langle (e _ {i}) _ {\mathrm{V}} v, (e _ {i}) _ {\mathrm{V}} v \rangle > 0, \end{array}
\]

故 \(a\) 成立。

令 B 为 \(t_{C}^{*}\) 上的逆型，它是限制到 \(t_{C}\) 的、由 F 经扩张标量从 \(g_{C}\) 上诱导的双线性型的逆型。根据第 VIII 章第 6 节第 4 号命题 7 的推论，有 \(\tilde{\Gamma}(\tau) = \mathrm{B}(\delta(\lambda), \delta(\lambda + 2\rho))\)。将 \(\delta : \mathrm{X(T)} \to t_{\mathbf{C}}^{*}\) 延拓为从 \(\mathrm{X(T)} \otimes \mathbf{R}\) 到 \(t_{C}^{*}\) 的 R-线性映射，并令 \(Q_{\Gamma}\) 为二次型 \(x \mapsto \mathrm{B}(\delta(x), \delta(x))\)，定义在 \(\mathrm{X(T)} \otimes \mathbf{R}\) 上；它是分离的且在 W 下不变，并且

\[
\tilde {\Gamma} (\tau) = \mathrm{B} (\delta (\lambda + \rho), \delta (\lambda + \rho)) - \mathrm{B} (\delta (\rho), \delta (\rho)) = \mathrm{Q} _ {\Gamma} (\lambda + \rho) - \mathrm{Q} _ {\Gamma} (\rho).
\]

我们证明型 \(Q_{\Gamma}\) 是正的。事实上，若 \(x \in \mathrm{X}(\mathrm{T}) \otimes \mathbf{Q}\)，则元素 \(\delta(x)\) 属于 \(t_{C}^{*}\)，它在 t 上取纯虚值，因而在 it 上取实值；注意到对于 \(y \in it\)，有 \(\mathrm{F}(y, y) \geq 0\)，即可得出结论。

还需证明 b) 中的唯一性断言。令 Q 为 \(X(T) \otimes \mathbf{R}\) 上满足所需条件的二次型，令 \(\Phi\)（分别为 \(\Phi_{\Gamma}\)）为与 Q（分别与 \(Q_{\Gamma}\)）关联的双线性型。对于 \(\lambda, \mu \in X_{++}\)，有

\[
\begin{array}{l} \Phi (\lambda , \mu) = (Q (\lambda + \mu + \rho) - Q (\rho)) - (Q (\lambda + \rho) - Q (\rho)) - (Q (\mu + \rho) - Q (\rho)) \\ \qquad = \Phi_ {\Gamma} (\lambda , \mu). \end{array}
\]

由于 \(\mathrm{X(T)_{++}}\) 生成 R-向量空间 \(\mathrm{X(T)}\otimes\mathbf{R}\)，有 \(\Phi=\Phi_{\Gamma}\)，所以 \(Q=Q_{\Gamma}\)。

注. 令 \(x \in g\)。存在严格正实数 A，使得对于每个不可约表示 \(\tau : G \to \mathbf{GL}(V)\) 以及 V 上每个在 G 下不变的 Hilbert 结构，

\[
\| \mathrm{L} (\tau) (x) \| ^ {2} \leq \mathrm{A}. \tilde {\Gamma} (\tau).
\]

事实上，沿用前面证明中的记号，可以选取 g 的基 \((e_{i})\)，使得 \(x = ae_{1}, a \in R\)。于是对于 \(v \in V\)，有

\[
\langle x _ {\mathrm{V}} v, x _ {\mathrm{V}} v \rangle = | a | ^ {2} \langle e _ {1} v, e _ {1} v \rangle \leq | a | ^ {2} \tilde {\Gamma} (\tau) \langle v, v \rangle .
\]

\section{Fourier 变换}

沿用上一段的记号和约定。